\documentclass[10pt]{article} % For LaTeX2e
%\usepackage{tmlr}
% If accepted, instead use the following line for the camera-ready submission:
%\usepackage[accepted]{tmlr}
% To de-anonymize and remove mentions to TMLR (for example for posting to preprint servers), instead use the following:
\usepackage[preprint]{tmlr}

% Optional math commands from https://github.com/goodfeli/dlbook_notation.
%\input{math_commands.tex}

\usepackage{hyperref}
\usepackage{url}
% amsmath, amssymb, natbib, graphicx, url, algorithm2e : Auto load
\usepackage[T1]{fontenc}
\usepackage[utf8]{inputenc}
\usepackage{amsmath}
\usepackage{amssymb}
\usepackage{amsthm}
\usepackage{tikz}
\usepackage{graphicx}
%\usepackage{pgfplots}
\usepackage{url}
\usepackage{xcolor}
\usepackage{bbm}
%\usepackage[ruled,vlined]{algorithm2e}
\usepackage{xfrac}
%\usepackage{subcaption}
%\usepackage[pagebackref=true]{hyperref}
\usepackage{mathtools}
%\usepackage[numbers]{natbib}
%\usepackage{url}
\usepackage{eufrak}
%\usepackage{pdfpages}
\usepackage{bbold}
%\usepackage{exscale,cmbright} %%%%%%%%%%%%-------- FONT
\usepackage{algorithm2e}
\usepackage{parskip}
%\RestyleAlgo{ruled}
\usepackage{aliascnt}
\usepackage{cleveref}

\DeclareMathOperator*{\argmax}{arg\,max}
\DeclareMathOperator*{\argmin}{arg\,min}


\Crefname{fact}{Fact}{Facts}
\Crefname{definition}{Definition}{Definitions}
\Crefname{lemma}{Lemma}{Lemmas}
\Crefname{corollary}{Corollary}{Corollarys}
\Crefname{example}{Example}{Examples}
\Crefname{remark}{Remark}{Remarks}
\Crefname{proposition}{Proposition}{Propositions}

\newtheorem{theorem}{Theorem}
\newtheorem{fact}{Fact}
\newtheorem{definition}{Definition}
\newtheorem{lemma}{Lemma}
\newtheorem{corollary}{Corollary}
\newtheorem{example}{Example}
\newtheorem{remark}{Remark}
\newtheorem{proposition}{Proposition}



% MACROS DIVERSES
\usepackage[utf8]{inputenc}
\usepackage{amsmath}
\usepackage{amssymb}
\usepackage{amsthm}
\usepackage{tikz}
\usepackage{graphicx}
%\usepackage{pgfplots}
\usepackage{url}
\usepackage{xcolor}
%\usepackage{bbm}
\usepackage{dsfont}
%\usepackage[ruled,vlined]{algorithm2e}
\usepackage{xfrac}
%\usepackage{subcaption}
\usepackage[]{hyperref}
\usepackage{mathtools}
%\usepackage[numbers]{natbib}
\usepackage{url}
%\usepackage{eufrak}
%\usepackage{pdfpages}
\usepackage{bbold}
%\usepackage{exscale,cmbright} %%%%%%%%%%%%-------- FONT
%\usepackage{algorithm2e}
\usepackage{parskip}
%\RestyleAlgo{ruled}
\usepackage{aliascnt}
\usepackage{cleveref}
\usepackage{subcaption}
\usepackage{makecell}
\usepackage{array}
% To allow \sout{} (strikeout)
\usepackage{ulem}
\normalem % to ensure \emph{} does not become ``underlined''
%
\usepackage{tikz}
\usepackage{svg}

\usepackage{makecell}

\usepackage[a4paper, total={6in, 8in}]{geometry}



\def\P{\mathbb{P}}
\def\one{\mathds{1}}
\def\Ind{\mathds{1}}

% MACROS DIVERSES
\newcommand{\R}[0]{\mathbb{R}}
\newcommand{\C}[0]{\mathbb{C}}
\newcommand{\N}[0]{\mathbb{N}}
\newcommand{\Z}[0]{\mathbb{Z}}
\newcommand{\Q}[0]{\mathbb{Q}}

\newcommand{\Prob}[0]{\mathbb{P}}
\newcommand{\prob}[0]{\mathbb{p}}
\newcommand{\E}[0]{\mathbb{E}}
\newcommand{\V}[0]{\mathbb{V}}
%\newcommand{\Ind}[0]{\mathbbm{1}}

\newcommand{\eqdef}{\vcentcolon=}

\newcommand\mech[1]{\mathfrak{#1}}
%\newcommand\mech[1]{\mathcal{#1}}
\newcommand\set[1]{\mathcal{#1}}
\newcommand\vect[1]{\mathbf{#1}}

\newcommand\dom[1]{\operatorname{dom}\left({#1}\right)}
\newcommand\codom[1]{\operatorname{codom}\left({#1}\right)}

\newcommand\ninf[0]{n \rightarrow +\infty}

\newcommand\conv[1]{\overset{}{\underset{#1}{\longrightarrow}}}
\newcommand\distconv[1]{\overset{\mathcal{L}}{\underset{#1}{\longrightarrow}}}
\newcommand\probconv[1]{\overset{\Prob}{\underset{#1}{\longrightarrow}}}
\newcommand\asconv[1]{\overset{\text{a.s.}}{\underset{#1}{\longrightarrow}}}

\newcommand{\iid}{\stackrel{\mathclap{\smallfont\text{i.i.d.}}}{\sim}}

\newcommand\tv[2]{\mathrm{TV}\left( {#1}, {#2} \right)}
\newcommand\kl[2]{\mathrm{KL}\left( \left. {#1}\right\Vert {#2} \right)}
\newcommand\renyi[3]{\text{D}_{#1}\left( \left. {#2}\right\Vert {#3} \right)}

\newcommand\ham[2]{d_\mathrm{ham}\left( {#1}, {#2} \right)}
\newcommand\hamext[1]{d_\mathrm{ham}\left( {#1}\right)}

\newcommand\intslice[2]{\left\{ {#1}, \dots,  {#2} \right\}}

\newcommand\p[1]{\left( {#1}\right)}
\newcommand\ceil[1]{\left\lceil {#1}\right\rceil}
\newcommand\card[1]{\# \left( {#1}\right)}
%
\newcommand\RG[1]{\textcolor{red}{#1}}
\newcommand\mRG[1]{\marginpar{\tiny\RG{#1}}}
\newcommand\AG[1]{\textcolor{magenta}{#1}}
\newcommand\CL[1]{\textcolor{blue}{#1}}

\setsvg{inkscapeexe=inkscape}
\setsvg{inkscapeopt=-z -D}


\title{About the Cost of Global Privacy in Density Estimation}

% Authors must not appear in the submitted version. They should be hidden
% as long as the tmlr package is used without the [accepted] or [preprint] options.
% Non-anonymous submissions will be rejected without review.

\author{{%
 \name Clément Lalanne \email clement.lalanne@ens-lyon.fr \\
 \addr Univ. Lyon, ENS Lyon, UCBL, CNRS, Inria,  LIP, F-69342, Lyon Cedex 07, France
 \AND
 \name Aurélien Garivier \email aurelien.garivier@ens-lyon.fr \\
 \addr Univ. Lyon, ENS Lyon, UMPA UMR 5669, 46 allée d'Italie, F-69364, Lyon cedex 07%
 \AND
 \name Rémi Gribonval \email remi.gribonval@inria.fr\\
 \addr Univ. Lyon, ENS Lyon, UCBL, CNRS, Inria,  LIP, F-69342, Lyon Cedex 07, France%
}}

% The \author macro works with any number of authors. Use \AND 
% to separate the names and addresses of multiple authors.

\newcommand{\fix}{\marginpar{FIX}}
\newcommand{\new}{\marginpar{NEW}}

\def\month{MM}  % Insert correct month for camera-ready version
\def\year{YYYY} % Insert correct year for camera-ready version
\def\openreview{\url{https://openreview.net/forum?id=XXXX}} % Insert correct link to OpenReview for camera-ready version


\begin{document}


\maketitle

\begin{abstract}%
  We study non-parametric density estimation for densities in Lipschitz and Sobolev spaces, and under \emph{global} privacy. In particular, we investigate regimes where the privacy budget is \emph{not} supposed to be constant. We consider the classical definition of global differential privacy, but also the more recent notion of global concentrated differential privacy. We recover the result of \citet{barber2014privacy} stating that histogram estimators are optimal against Lipschitz distributions for the $L^2$ risk and, under regular differential privacy, we extend it to other norms and notions of privacy. Then, we investigate higher degrees of smoothness, drawing two conclusions: First, and contrary to what happens with \emph{constant} privacy budget \citep{wasserman2010statistical}, there \emph{are} regimes where imposing privacy degrades the regular minimax risk of estimation on Sobolev densities. Second, so-called projection estimators are near-optimal against the same classes of densities in this new setup with pure differential privacy, but contrary to the constant privacy budget case, it comes at the cost of relaxation. With zero concentrated differential privacy, there is no need for relaxation, and we prove that the estimation is optimal.
\end{abstract}

%\CL{IMPORTANT : Avant de poursuivre la relecture, attendez ! J'ai trouvé une preuve pour les bornes inf sous $\rho$-zCDP qui donne des résultats qui matchent les bornes sup !! Je vais ajouter un paragraphe sur la methode de Assouad sous $\rho$-zCDP et adapter les preuves et théorèmes !}

\section{Introduction}

The communication of information built on users' data leads to new challenges, and notably privacy concerns. It is now well documented that the release of various quantities can, without further caution, have disastrous repercussions \citep{narayanan2006break,backstrom2007wherefore,fredrikson2015model,dinur2003revealing,homer2008resolving,loukides2010disclosure,narayanan2008robust,sweeney2000simple,gonon2023sparsity,wagner2018technical,sweeney2002k}.
In order to address this issue, differential privacy (DP) \citep{dwork2006calibrating} has become the gold standard in privacy protection. The idea is to add a proper layer of randomness %(i.e. the estimator does not only build on $X_1, \dots, X_n$ but also on another source of randomness) 
in order to hide each user's data. It is notably used by the US Census Bureau~\citep{abowd2018us}, Google~\citep{erlingsson2014rappor}, Apple~\citep{thakurta2017learning} and Microsoft~\citep{ding2017collecting}, among many others. 

As with other forms of communication or processing constraints \citep{BarnesHO19,BarnesHO20,AcharyaCST21a,AcharyaCMT21,AcharyaCST21b,AcharyaCFST21}, privacy recently gained a lot of attention from the statistical and theoretical machine learning communities. At this point, the list of interesting publications is far too vast to be exhaustive, but here is a sample : \citet{wasserman2010statistical} is the first article to consider problems analogous to the ones presented in this article. It notably studies the problem of nonparametric density estimation, to which we provide many complements. \citet{duchi2013localpreprint,duchi2013local,duchi2018minimax,barber2014privacy,acharya2021differentially,lalanne2022statistical} present general frameworks for deriving minimax lower-bounds under privacy constraints. Many parametric problems have already been studied, notably in \citet{acharya2018differentially,acharya2021differentially,karwa2017finite,kamath2019highdimensional,biswas2020coinpress,lalanne2022private,lalanne2023private,kamath2022improved,singhal2023polynomial}. Recently, some important contributions were made. For instance, \citet{asi2023robustness} sharply characterized the equivalence between private estimation and robust estimation with the inverse sensitivity mechanism \citep{asi2020near,asi2021nips}, and \citet{kamath2023biasvarianceprivacy} detailed the bias-variance-privacy trilemma
, proving in particular the necessity of adding bias, even on distributions with bounded support, for many private estimation problems.

We address here the problem of privately estimating a probability density, which fits in this line of work. Given $\vect{X} \eqdef \p{X_1, \dots, X_n} \sim \Prob_\pi^{\otimes n}$, where $\Prob_\pi$ refers to a distribution of probability that 
has a density $\pi$ 
%that is absolutely continuous 
with respect to the Lebesgue measure on $[0, 1]$, how to estimate $\pi$ privately? Technically, what metrics or hypothesis should be set on $\pi$? What is the cost of privacy? Are the methods known so far optimal? Such are the questions that are investigated in the rest of this article.

\subsection{Related work}

Non-parametric density estimation has been an important topic of research in statistics for many decades now. Among the vast literature on the topic, let us just mention the important references \citet{gyorfi2002distribution,tsybakov2003introduction}. 

Recently, the interest for private statistics has shone a new light on this problem. Remarkable early contributions \citep{wasserman2010statistical,hall2013differential} adapted histogram estimators, so-called projection estimators and kernel estimators to satisfy the privacy constraint. They conclude that the minimax rate of convergence, $n^{- 2 \beta/(2 \beta + 1)}$, where $n$ is the sample size and $\beta$ is the (Sobolev) smoothness of the density, is not affected by \emph{global} privacy. However, an important implicit hypothesis in this line of work is that $\epsilon$, the parameter that decides how private the estimation needs to be, is supposed not to depend on the sample size. This hypothesis may seem disputable, and more importantly, it fails to precisely characterize the tradeoff between utility and privacy. To the best of our knowledge, the only piece of work that studies this problem under \emph{global} privacy when $\epsilon$ is not supposed constant is \citet{barber2014privacy}. They study histogram estimators on Lipschitz distributions for the integrated risk. They conclude that the minimax risk of estimation is $\max \p{n^{-2/3} + (n\epsilon)^{-1}}$, showing how small $\epsilon$ can be before the minimax risk of estimation is degraded. Our article extends such results to high degrees of smoothness, to other definitions of \emph{global} differential privacy, and to other risks.

The literature under the much stricter notion of \emph{local} privacy is a lot richer. Contrary to \emph{global} privacy, \emph{local} privacy requires that each data holder anonymizes its data before them being communicated to an aggregator. It is a stronger definition of privacy than global differential privacy. A remarkable early piece of work \cite{duchi2018minimax} has brought a nice toolbox for deriving minimax lower bounds under local privacy that has proven to give sharp results for many problems. As a result, the problem of non-parametric density estimation (or its analogous problem of non-parametric regression) has been extensively studied under local privacy. For instance, \citet{butucea2020local} investigates the elbow effect and questions of adaptivity over Besov ellipsoids. \citet{kroll2021density} and \citet{schluttenhofer2022adaptive} study the density estimation problem at a given point with an emphasis on adaptivity. Universal consistency properties have recently been derived in \citet{gyorfi2023multivariate}. Analogous regression problems have been studied in \citet{berrett2021strongly} and in \citet{gyorfi2022rate}. Finally, the problem of optimal non-parametric testing has been studied in \citet{lam2022minimax}.

%More recently, similar lower-bounding toolboxes have been designed for global privacy \cite{acharya2021differentially,kamath2022improved,lalanne2022statistical}. This work builds on these toolboxes in order to study the problem of non-parametric density estimation under global privacy, when $\epsilon$ is not constant, and for a general smoothness parameter $\beta$.

\subsection{Contributions}

In this article, we investigate the impact of \emph{global} privacy when the privacy budget is not constant. We treat multiple definitions of \emph{global} privacy and different levels of smoothness for the densities of interest. 

In terms of upper-bounds, we analyze histogram and projection estimators at a resolution that captures the impact of the privacy and smoothness parameters. We also prove new lower bounds using the classical packing method combined with new tools that characterize the testing difficulty under global privacy from \citet{acharya2021differentially,kamath2022improved,lalanne2022statistical}.

In particular, for Lipschitz densities and under pure differential privacy, we recover the results of \citet{barber2014privacy} with a few complements. 
We then extend the estimation on this class of distributions to the context of concentrated differential privacy \citep{bun2016concentrated}, a more modern definition of privacy that is compatible with stochastic processes relying on Gaussian noise. We finally investigate higher degrees of smoothness by looking at periodic Sobolev distributions.
The main results are summarized in \Cref{table:mainresults}.

\begin{table}[ht]
\centering
\label[table]{table:mainresults}
\makebox[\textwidth][c]{
\begin{tabular}{ |c||c|c| } 
 \hline
        & \textbf{$\epsilon$-DP} \Cref{eq:defdp} & \textbf{$\rho$-zCDP} \Cref{eq:defcdp} \\ 
 \hline\hline
 \makecell{\textbf{Lipschitz} \\ \Cref{eq:deflip}} & 
 \makecell{
 \underline{Upper-bound: } \\
 $O\p{\max \left\{ n^{-2/3}, (n\epsilon)^{-1} \right\}}$\\
 \citep{barber2014privacy} \& \Cref{proposition:histdp} \vspace{0.2cm}\\
 \hline
 \underline{Lower-bounds: } \\
      -Pointwise: $\Omega\p{\max \left\{ n^{-2/3}, (n\epsilon)^{-1} \right\}}$ \\
     \Cref{th:lpdp} \& \Cref{cor:linfdp}\\
      -Integrated: $\Omega\p{\max \left\{ n^{-2/3}, (n\epsilon)^{-1} \right\}}$ \\
     \citep{barber2014privacy} \& \Cref{th:lidp}
 } & 
 \makecell{
 \underline{Upper-bound: } \\
 $O\p{\max \left\{ n^{-2/3}, (n\sqrt{\rho})^{-1} \right\}}$\\
 \Cref{proposition:histdp} \vspace{0.2cm}\\
 \hline
 \underline{Lower-bounds: } \\
      -Pointwise: $\Omega\p{\max \left\{ n^{-2/3}, (n\sqrt{\rho})^{-1} \right\}}$ \\
     \Cref{th:lpdp} \& \Cref{cor:linfdp}\\
      -Integrated: $\Omega\p{\max \left\{ n^{-2/3}, (n\sqrt{\rho})^{-1} \right\}}$ \\
     \Cref{th:lidp}
 } \\ 
 \hline
 \makecell{\textbf{Periodic Sobolev} \\ Smoothness $\beta$ \\ \Cref{eq:defpsob}} & 
 \makecell{
 \underline{Upper-bounds: } \\
-Pure DP: $O\p{\max \left\{ n^{-\frac{2 \beta}{2 \beta + 1}}, (n \epsilon)^{-\frac{2 \beta}{ \beta + 3/2}} \right\}}$ \\
     \Cref{proposition:projdp}\\
-Relaxed: $\tilde{O}\p{\max \left\{ n^{-\frac{2 \beta}{2 \beta + 1}}, (n \epsilon)^{-\frac{2 \beta}{ \beta + 1}} \right\}} $ \\
     $\tilde{O}$ hides polylog factors. \Cref{proposition:relaxation} \vspace{0.2cm}\\
     \hline
 \underline{Lower-bound: } \\
      $\Omega\p{\max \left\{ n^{-\frac{2 \beta}{2 \beta + 1}}, (n \epsilon)^{-\frac{2 \beta}{ \beta + 1}} \right\}}$ \\
     \Cref{th:sidp}
 } & 
 \makecell{
 \underline{Upper-bound: } \\
$O\p{\max \left\{ n^{-\frac{2 \beta}{2 \beta + 1}}, (n \sqrt{\rho})^{-\frac{2 \beta}{ \beta + 1}} \right\}}$ \\
     \Cref{proposition:projdp}\vspace{0.2cm}\\
     \hline
 \underline{Lower-bound: } \\
      $\Omega\p{\max \left\{ n^{-\frac{2 \beta}{2 \beta + 1}}, (n \sqrt{\rho})^{-\frac{2 \beta}{ \beta + 1}} \right\}}$ \\
     \Cref{th:sidp}
 } \\ 
 \hline
\end{tabular}
}
\caption{Summary of the results}
\end{table}

%For the class of Lipschitz distributions, we recover the known result that histogram estimators are optimal against $\epsilon$-DP mechanisms, and we show that this result also holds for the pointwise risk (and as a consequence in squared infinite norm as well). Against $\rho$-zCDP mechanisms on the other hand, the optimality raises new questions. We show that in this case, histogram estimators are optimal for the pointwise risk, and hence also in squared infinite norm. Yet, we have a small gap between our upper and lower-bounds for the integrated risk. 

%For the class of Sobolev distributions, we again have a small gap between upper-bounds and lower-bounds for both definitions of privacy. Since in the case $\beta = 1$ (the smoothness parameter), one usually expects to recover the same results as with Lipschitz densities \cite{tsybakov2003introduction}, we suspect that our lower-bounds are correct in the setup of $\epsilon$-DP, but that the mechanism realizing the upper-bound is either suboptimal, either suboptimally analyzed. However, via a simple relaxation to $(\epsilon, \delta)$-differential privacy with polynomially small $\delta$, we show that it is possible to obtain a mechanism that matches the lower-bound up to polylog factors. For concentrated differential privacy, we do not know if the small gaps come from the upper or lower-bounds, or from both.

%In both cases, the gaps are small, especially when the smoothness of the target density is high. This piece of work demonstrates that there is still room for improvement on the topic of density estimation under global privacy.

The paper is organized as follows. 
The required notions regarding \emph{global} differential privacy are recalled in \Cref{sec:reminders}. Histogram estimators and projection estimators are respectively studied in \Cref{sec:histogram}, on Lipschitz densities, and in \Cref{sec:projection}, on periodic Sobolev densities. A short conclusion in provided in Section~\ref{sec:ccl}.

\section{Global Differential Privacy}
\label{sec:reminders}
We recall in this section some useful notions of \emph{global} privacy.
Here, $\set{X}$ and $n$ refer respectively to the sample space and to the sample size.
Given $\epsilon > 0$ and $\delta \in [0, 1)$, a randomized mechanism $\mech{M} : \set{X}^n \rightarrow \codom{\mech{M}}$ (for \emph{codomain} or image of $\mech{M}$) is $(\epsilon, \delta)$-differentially private (or $(\epsilon, \delta)$-DP) \citep{dwork2006calibrating,dwork2006our} if for all $\vect{X}, \vect{Y} \in \set{X}^n$ and all measurable $S \subseteq \codom{\mech{M}}$:
\begin{equation}
\label{eq:defdp}
    \ham{\vect{X}}{\vect{Y}} \leq 1 \implies \Prob_{\mech{M}}\p{\mech{M}(\vect{X}) \in S} \leq e^\epsilon \Prob_{\mech{M}}\p{\mech{M}(\vect{Y}) \in S} + \delta \;,
\end{equation}
where $\ham{\cdot}{\cdot}$ denotes the Hamming distance on $\set{X}^n$. 

In order to sharply count the privacy of a composition of many Gaussian mechanisms (see~\citet{abadi2016deep}), privacy is also often characterized in terms of Renyi divergence \citep{mironov2017renyi}. Nowadays, it seems that all these notions tend to converge towards the definition of \emph{zero concentrated differential privacy} \citep{dwork2016concentrated,bun2016concentrated}. Given $\rho \in (0, +\infty)$, a randomized mechanism $\mech{M} : \set{X}^n \rightarrow \codom{\mech{M}}$ is $\rho$-zero concentrated differentially private ($\rho$-zCDP) if for all $\vect{X}, \vect{Y} \in \set{X}^n$, %for all $\alpha \in (1, \infty)$,
\begin{equation}
\label{eq:defcdp}
    \ham{\vect{X}}{\vect{Y}} \leq 1 \implies   \forall 1 < \alpha < +\infty, \renyi{\alpha}{\mech{M}(\vect{X})}{\mech{M}(\vect{Y})} \leq \rho \alpha 
\end{equation}
where $\renyi{\alpha}{\cdot}{\cdot}$ denotes the Renyi divergence of level $\alpha$, defined when $\alpha>1$ as 
\begin{equation*}
    \renyi{\alpha}{\mathbb{P}}{\mathbb{Q}} \eqdef \frac{1}{\alpha - 1} \ln \int  \p{\frac{d \mathbb{P}}{ d \mathbb{Q}}}^{\alpha - 1} d \mathbb{Q} \;.
\end{equation*}
For more details, we recommend referring to the excellent article \citet{van2014renyi}.

There are links between $(\epsilon, \delta)$-DP and $\rho$-zCDP. For instance, if a mechanism is $\rho$-zCDP, then  \citep[Proposition 3]{bun2016concentrated} it is $(\epsilon, \delta)$-DP for a collection of $(\epsilon, \delta)$'s that depends on $\rho$. Conversely,  if a mechanism is $(\epsilon, 0)$-DP, then \citep[Proposition 4]{bun2016concentrated} it is also $\epsilon^2/2$-zCDP.


Given a deterministic function $f$ mapping a dataset to a quantity in $\R^d$, the Laplace mechanism \citep{dwork2006calibrating} and Gaussian mechamisn \citep{bun2016concentrated} are two famous ways to turn $f$ into a private mechanism.
Defining the $l_1$ sensitivity of $f$ as
\begin{equation*}
\begin{aligned}
    \Delta_1 f \eqdef \sup_{\vect{X}, \vect{Y} \in \set{X}^n : \ham{\vect{X}}{\vect{Y}} \leq 1} \big\| f (\vect{X})  - f (\vect{Y})\big\|_1 \;,
\end{aligned}
\end{equation*}
the Laplace mechanism instantiated with $f$ and $\epsilon > 0$ is defined as
\begin{equation*}
    \begin{aligned}
        \vect{X} \mapsto f(\vect{X}) + \frac{\Delta_1 f}{\epsilon} \mathcal{L}(I_d) \;,
    \end{aligned}
\end{equation*}
where $\mathcal{L}(I_d)$ refers to a random vector of dimension $d$ with independent components that follow a centered Laplace distribution of parameter $1$. 
Notice that we took the liberty to use the same notation for the random variable and for its distribution. We made this choice for brevity, and because it does not really create any ambiguity.
It is $(\epsilon, 0)$-DP (simply noted $\epsilon$-DP) \citep{dwork2006our,dwork2006calibrating}.
Likewise, defining the $L^2$ sensitivity of $f$ as
\begin{equation*}
\begin{aligned}
    \Delta_2 f \eqdef \sup_{\vect{X}, \vect{Y} \in \set{X}^n : \ham{\vect{X}}{\vect{Y}} \leq 1} \big\| f (\vect{X})  - f (\vect{Y})\big\|_2 \;,
\end{aligned}
\end{equation*}
the Gaussian mechanism instantiated with $f$ and $\rho > 0$ is defined as
\begin{equation*}
    \begin{aligned}
        \vect{X} \mapsto f(\vect{X}) + \frac{\Delta_2 f}{\sqrt{2 \rho}} \mathcal{N}(0, I_d) \;,
    \end{aligned}
\end{equation*}
where $\mathcal{N}(0, I_d)$ refers to a random vector of dimension $d$ with independent components that follow a centered Normal distribution of variance $1$. It is $\rho$-zCDP \citep{bun2016concentrated}.
%\AG{dans ces définitions on confond un peu la notation de la loi d'un vecteur aléatoire, avec le vecteur lui-même; bon, c'est peut-être pas grave.}

\paragraph{A quick word on local privacy.}

Global privacy comes with the hypothesis of a trusted aggregator that sees the entire dataset, and builds an estimator with it. Only the produced estimator is private. In order to give an example, this is like having a datacenter that stores all the information about the users of a service, but only outputs privatized statistics.

Local privacy on the other hand does not make that hypothesis. Each piece of data is anonymized locally (on the user's device) and then it is communicated to an aggregator. Any locally private mechanism is also globally private, but the converse is not true.

At first, local privacy can seem more appealing : it is indeed a stronger notion of privacy. However, it degrades the utility \emph{much more} than global privacy. As a result, both notions are interesting, and the use of one or the other must be weighted for a given problem. This work focuses on the \emph{global} variant.

\section{Histogram Estimators and Lipschitz Densities}
\label{sec:histogram}


Histogram estimators approximate densities with a piecewise continuous function by counting the number of points that fall into each bin of a partition of the support. Since those numbers follow binomial distributions, the study of histogram estimators is rather simple. Besides, they are particularly interesting when privacy is required, since the sensitivity % (see \Cref{sec:reminders}) 
of a histogram query is bounded independently of the number of bins. They were first studied in this setup in \citet{wasserman2010statistical}, while \citet{barber2014privacy} provided new lower-bounds that did not require a constant privacy budget.

As a warm-up, this section proposes a new derivation of known results in more modern lower-bounding frameworks \citep{acharya2021differentially,kamath2022improved,lalanne2022statistical}, and then extends these upper-bounds and lower-bounds to the case of zCDP. Furthermore, it also covers the pointwise risk as well as the infinite-norm risk.

Let $h>0$ be a given bandwidth or binsize. In order to simplify the notation, we suppose  without loss of generality that $1/h \in \N \setminus \{0 \}$ (if the converse is true, simply take $h' = 1/\ceil{1/h}$ where $\ceil{x}$ refers to the smallest integer bigger than $x$). 
$[0, 1]$ is partitioned in $\frac{1}{h}$ sub-intervals of length $h$, which are called the bins of the histogram.
Let $Z_1, \dots, Z_{1/h}$ be independent and identically distributed random variables with the same distribution as a random variable $Z$ that is supposed to be centered and to have a finite variance. Given a dataset $\vect{X} = (X_1, \dots, X_n)$, the (randomized) histogram estimator is defined for $x\in[0,1]$ as 
\begin{equation}
\label{eq:hist}
    \hat{\pi}^{\mathrm{hist}} (\vect{X}) (x) \eqdef \sum_{b \in \mathrm{bins}} \Ind_b(x) \frac{1}{nh} \p{\sum_{i=1}^n \Ind_b(X_i) +Z_b}\;.
\end{equation}
We indexed the $Z$'s by a bin instead of an integer without ambiguity.
Note that by taking $Z$ almost-surely constant to $0$, one recovers the usual (non-private) histogram estimator of a density. % In the sequel, $Z$ is used to obtain privacy.


\subsection{General utility of histogram estimators}

Characterizing the utility of \eqref{eq:hist} typically requires assumptions on the distribution $\pi$ to estimate. The class of $L$-Lipschitz densities is defined as
\begin{equation}
\label{eq:deflip}
\Theta^{\mathrm{Lip}}_L \eqdef \left\{ \pi \in \mathcal{C}^0([0, 1], \R_+) \left| 
\begin{cases}
  \forall x, y \in [0, 1], |\pi(y) - \pi(x)| \leq L |y-x|\;, \\
  \int_{[0, 1]} \pi = 1 \;.  
\end{cases}  
 \right.\right\} \;.
\end{equation}
The following general-purpose lemma gives an upper-bound on the error that the histogram estimator makes on Lipschitz distributions:
\begin{lemma}[General utility of \eqref{eq:hist}]
\label[lemma]{lemma:histogramutility}
There exists $C_L > 0$, a positive constant that only depends on $L$, such that
\begin{equation*}
    \sup_{x_0 \in [0, 1]} \sup_{\pi \in \Theta^{\mathrm{Lip}}_L}
    \E_{\vect{X} \sim \Prob_{\pi}^{\otimes n}, \hat{\pi}^{\mathrm{hist}}}\p{ \p{\hat{\pi}^{\mathrm{hist}}(\vect{X})(x_0) - \pi(x_0)}^2}
    \leq C_L \p{h^2 + \frac{1}{nh} + \frac{\V(Z)}{n^2 h^2}} \;.
\end{equation*}
\end{lemma}
The proof is given in~\Cref{proofoflemma:histogramutility}.
The term $h^2$ corresponds to the bias of the estimator. The variance term $\frac{1}{nh} + \frac{\V(Z)}{n^2 h^2}$ exhibits two distinct contributions : the sampling noise $\frac{1}{nh}$ and the privacy noise $\frac{\V(Z)}{n^2 h^2}$. In particular, the utility of $\hat{\pi}^{\mathrm{hist}}$ changes depending whether the variance is dominated by the sampling noise or by the privacy noise.

\subsection{Privacy and bin size tuning}

$\hat{\pi}^{\mathrm{hist}} (\vect{X})$ is a simple function of the bin count vector
$f(\vect{X}) \eqdef \p{\sum_{i=1}^n \Ind_{b_1}(X_i), \dots,  \sum_{i=1}^n \Ind_{b_{1/h}}(X_i)}$. In particular, since the bins form a partition of $[0, 1]$, changing the value of one of the $X$'s can change the values of at most two components of $f(\vect{X})$ by at most $1$. Hence, the $l_1$ and $l_2$ sensitivities of $f$ are respectively $2$ and $\sqrt{2}$. By a direct application of the Laplace or Gaussian mechanisms, and by choosing the binsize that minimizes the variance, we obtain the following privacy-utility result :
\begin{theorem}[Privacy and utility of \eqref{eq:hist} - DP case] 
\label[theorem]{proposition:histdp}
Given $\epsilon > 0$,
using $\hat{\pi}^{\mathrm{hist}}$ with $h =  \max( n^{-1/3},\\ (n\epsilon)^{-1/2} )$ and $Z = \frac{2}{\epsilon} \mathcal{L}(1)$, where $\mathcal{L}(1)$ refers to a random variable following a Laplace distribution of parameter $1$, leads to an $\epsilon$-DP procedure. Furthermore, in this case, there exists $C_L > 0$, a positive constant that only depends on $L$, such that
\begin{equation*}
    \sup_{x_0 \in [0, 1]} \sup_{\pi \in \Theta^{\mathrm{Lip}}_L}
    \E_{\vect{X} \sim \Prob_{\pi}^{\otimes n}, \hat{\pi}^{\mathrm{hist}}}\p{ \p{\hat{\pi}^{\mathrm{hist}}(\vect{X})(x_0) - \pi(x_0)}^2}
    \leq C_L  \max \left\{ n^{-2/3}, (n\epsilon)^{-1} \right\} \;.
\end{equation*}
Furthermore, given $\rho > 0$,
using $\hat{\pi}^{\text{hist}}$ with $h =  \max( n^{-1/3},(n\sqrt{\rho})^{-1/2} )$ and $Z = \sqrt{\frac{1}{\rho}} \mathcal{N}(0, 1)$, where $\mathcal{N}(0, 1)$ refers to a random variable following a centered Gaussian distribution of variance $1$, leads to a $\rho$-zCDP procedure. Furthermore, in this case, there exists $C_L > 0$, a positive constant that only depends on $L$, such that
\begin{equation*}
    \sup_{x_0 \in [0, 1]} \sup_{\pi \in \Theta^{\text{Lip}}_L}
    \E_{\vect{X} \sim \Prob_{\pi}^{\otimes n}, \hat{\pi}^{\text{hist}}}\p{ \p{\hat{\pi}^{\text{hist}}(\vect{X})(x_0) - \pi(x_0)}^2}
    \leq C_L  \max \left\{ n^{-2/3}, (n \sqrt{\rho})^{-1} \right\} \;.
\end{equation*}
\end{theorem}
Note that this bound is uniform in $x$. In particular, by integration on $[0, 1]$, the same bound also holds for the integrated risk (in $L^2$ norm). As expected, the optimal bin size $h$ depends on the sample size $n$ and on the parameter ($\epsilon$ or $\rho$) tuning the privacy.

\subsection{Lower-bounds and minimax optimality}

All lower-bounds will be investigated in a minimax sense. Given a class $\Pi$ of admissible densities, a semi-norm $\| \cdot \|$ on a space containing the class $\Pi$, and a non-decreasing positive function $\Phi$ such that $\Phi(0) = 0$, the minimax risk is defined as
\begin{equation*}
    \inf_{\hat{\pi} \text{ s.t. } \mathcal{C}} \sup_{\pi \in \Pi}  \E_{\vect{X} \sim \Prob_{\pi}^{\otimes n}, \hat{\pi}} \Phi(\|\hat{\pi}(\vect{X}) - \pi \|) \;,
\end{equation*}
where $\mathcal{C}$ is a condition that must satisfy the estimator (privacy in our case).

\paragraph{General framework.}
A usual technique for the derivation of minimax lower bounds on the risk uses a reduction to a testing problem (see \citet{tsybakov2003introduction}). Indeed, if a family $\Pi' \eqdef \{ \pi_1, \dots, \pi_m\} \subset \Pi$ of cardinal $m$ is an $\Omega$-packing of %$\Pi'$ (that is if $\Pi$ 
$\Pi$ (that is
if $i \neq j \implies \|\pi_i - \pi_j \| \geq 2 \Omega$), then
%\mRG{L'équation ci-dessous est un peu dure à digérer pour un lecteur qui ne l'a jamais vue je pense \CL{Vu qu'on donne les refs, je pense que c'est bon}}
a lower bound is given by %\mRG{Dans cette equation, ne faudrait-il pas a) remplacer $\mech{M}$ par $\hat{\pi}$; b) mettre $\vect{X} \sim \Prob_{\pi}^{\otimes n}$ au lieu de $\vect{X} \sim \Prob_{\pi}$ a droite; c) remplacer $\sup$ par $\max$?}
\begin{equation}
\label{jkhgqskdjhfbgkjqhwgsdkfjhgkqwsjhdgfjkhwsgdb}
\begin{aligned}
    \inf_{\hat{\pi} \text{ s.t. } \mathcal{C}} \sup_{\pi \in \Pi}  \E_{\vect{X} \sim \Prob_{\pi}^{\otimes n}, \hat{\pi}} &\Phi(\|\hat{\pi}(\vect{X}) - \pi \|) \\
    &\geq 
    \Phi(\Omega) \underset{\Psi : \codom{\hat{\pi}} \rightarrow\intslice{1}{m}}{\inf_{\hat{\pi} \text{ s.t. } \mathcal{C}}}  \max_{i \in \intslice{1}{m}} \Prob_{\mathbf{X} \sim \Prob_{\pi_i}^{\otimes n},\hat{\pi}} \p{ \Psi \p{\hat{\pi}(\vect{X})} \neq i}\;.
\end{aligned}
\end{equation}
For more details, see \citet{duchi2018minimax,acharya2021differentially,lalanne2022statistical}. The right-hand side characterizes the difficulty of discriminating the distributions of the packing by a statistical test. Independently on the condition $\mathcal{C}$, it can be lower-bounded using information-theoretic results such a Le Cam's lemma \citep[Lemma 5.3]{rigollet2015high} or Fano's lemma \citep[Theorem 3.1]{giraud2021introduction}. When $\mathcal{C}$ is a local privacy condition, \citet{duchi2018minimax} provides analogous results that take privacy into account. Recent work \citep{acharya2021differentially,kamath2022improved,lalanne2022statistical} provides analogous forms for multiple notions of \emph{global} privacy. When using this technique, finding good lower-bounds on the minimax risk boils down to finding a packing of densities that are far enough from one another without being too easy to discriminate with a statistical test.

It is interesting to note that for the considered problem, this technique does not yield satisfying lower-bounds with $\rho$-zCDP every time Fano's lemma is involved. Systematically, a small order is lost. To circumvent that difficulty, we had to adapt Assouad's technique to the context of $\rho$-zCDP. Similar ideas have been used in \citet{duchi2018minimax} for lower-bounds under \emph{local} differential privacy and in \citet{acharya2021differentially} for regular \emph{global} differential privacy. To the best of our knowledge, such a technique has never been used in the context of global \emph{concentrated} differential privacy, and is presented in \Cref{assouadcdp}. In all the proofs of the lower-bounds, we systematically presented both approaches whenever there is a quantitative difference. This difference could be due to small suboptimalities in Fano's lemma for concentrated differential privacy, or simply to the use of a suboptimal packing.


\subsubsection{Pointwise lower-bound}

The first lower-bound that will be investigated is with respect to the pointwise risk. Pointwise, that is to say given
$x_0 \in [0, 1]$, the performance of the estimator $\hat{\pi}$ is measured by how well it approximates $\pi$ at $x_{0}$ with the quadratic risk $\E_{\vect{X} \sim \Prob_{\pi}^{\otimes n}, \hat{\pi}}\p{\p{\hat{\pi}(\vect{X})(x_0) - \pi(x_0)}^2}$.
%
Technically, it is the easiest since it requires a "packing" of only two elements, which gives the following lower-bound:

\begin{theorem}[Pointwise lower-bound]
\label{th:lpdp}
There exists $C_L>0$, %\mRG{Est-ce la meme $C_L$ que dans les précédents résultats ? \CL{Non mais il est possible de prendre la plus grande des deux} Je voulais dire: veut-on juste parler d'existence, ou donner l'expression ou au moins le comportement en $L$ ?}
a positive constant depending only on $L$ such that, for any $x_0 \in [0, 1]$, there exist $n_0(x_0, L) \in \N$ and $c_0(x_0, L) > 0$ such that for any $n \geq n_0$, and any $\alpha \geq c_0/n$
%$\epsilon$ such that $n \alpha \geq c_0$,
\begin{equation}
\label{eq:deftriangle}
    \inf_{\hat{\pi} \mathrm{ s.t. } \mathcal{C}} \sup_{\pi \in \Theta^{\mathrm{Lip}}_L}
    \E_{\vect{X} \sim \Prob_{\pi}^{\otimes n}, \hat{\pi}}\p{\p{\hat{\pi}(\vect{X})(x_0) - \pi(x_0)}^2} \geq C_L^{-1} \max \left\{ n^{-2/3}, (n\alpha)^{-1} \right\}
    \;,
\end{equation}
where $\alpha = \epsilon$ when the condition $\mathcal{C}$ is the $\epsilon$-DP condition and $\alpha = \sqrt{\rho}$ when %the condition 
$\mathcal{C}$ is $\rho$-zCDP.
\end{theorem}
\begin{proof}[Proof idea] %\AG{c'est bcp dire sketch of proof, c'est plutôt une déclaration d'intention. je ne la mettrais pas dans un environnement "proof"}
Let $x_0 \in [0, 1]$.
As explained above, finding a "good" lower-bound can be done by finding and analyzing a "good" packing of the parameter space. Namely, in this case, we have to find distributions on $[0, 1]$ that have a $L$-Lipschitz density (w.r.t. Lebesque's measure) such that the densities are far from one another at $x_0$, but such that it is not extremely easy to discriminate them with a statistical test.
We propose to use a packing $\{ \Prob_f, \Prob_g\}$ of two elements where $g$ is the constant function on $[0, 1]$ (hence $\Prob_g$ is the uniform distribution) and $f$ deviates from $g$ by a small triangle centered at $x_0$. 
%The two densities are represented in \Cref{fig:onetriangle}.

After analyzing various quantities about these densities, such as their distance at $x_0$, their KL divergences or their TV distance, we leverage Le Cam-type results to conclude.
\end{proof}
The full proof can be found in \Cref{proofofth:lpdp}.

%By comparing this result with \Cref{proposition:histdp}, it means that the histogram estimators has minimax optimal rates of convergence for the pointwise risk, and for both notions of global privacy.

Additionally, we can notice that, when applied to any fixed $x_0 \in [0, 1]$, \Cref{th:lpdp} immediately gives the following corollary for the control in infinite norm :

\begin{corollary}[Infinite norm lower-bound]
\label{cor:linfdp}
There exists $C_L>0$, a positive constant depending only on $L$ such that there exist $n_0(L) \in \N$ and $c_0(L) > 0$ such that for any $n \geq n_0$, and any 
$\alpha \geq c_0/n$
%$\epsilon$ such that $n \alpha \geq c_0$,
\begin{equation}
\label{eq:deftriangle}
    \inf_{\hat{\pi} \mathrm{ s.t. } \mathcal{C}} \sup_{\pi \in \Theta^{\mathrm{Lip}}_L}
    \E_{\vect{X} \sim \Prob_{\pi}^{\otimes n}, \hat{\pi}}\|\hat{\pi}(\vect{X})- \pi\|_{\infty}^2  \geq C_L^{-1} \max \left\{ n^{-2/3}, (n\alpha)^{-1} \right\}
    \;,
\end{equation}
where $\alpha = \epsilon$ when the condition $\mathcal{C}$ is the $\epsilon$-DP condition and $\alpha = \sqrt{\rho}$ when %the condition 
$\mathcal{C}$ is $\rho$-zCDP.
\end{corollary}

\paragraph{On the optimality and on the cost of privacy.}

\Cref{proposition:histdp}, \Cref{th:lpdp} and \Cref{cor:linfdp} give the following general result : Under $\epsilon$-DP or under $\rho$-zCDP, histogram estimators have minimax-optimal rates of convergence against distributions with Lipschitz densities, for the pointwise risk or the risk in infinite norm. In particular, in the \emph{low privacy} regime (``large'' $\alpha$), the usual minimax rate of estimation of $n^{-\frac{2}{3}}$ is not degraded. This includes the early observations of \citet{wasserman2010statistical} in the case of constant $\alpha$ ($\epsilon$ or $\sqrt{\rho}$). However, in the \emph{high privacy} regimes ($\alpha \ll n^{-\frac{1}{3}}$), these results prove a systematic degradation of the estimation. Those regimes are the same as in \citet{barber2014privacy}, the metrics on the other hand are different. % \mRG{Est-ce entierement nouveau ? A positionner par rapport à \cite{barber2014privacy} }

\subsubsection{Integrated lower-bound} 

The lower-bound of \Cref{th:lpdp} is interesting, but its pointwise (or in infinite norm in the case of \Cref{cor:linfdp}) nature means that much global information is possibly lost. 
Instead,
one can look at the integrated risk $\E_{\vect{X} \sim \Prob_{\pi}^{\otimes n}, \hat{\pi}}\|\hat{\pi}(\vect{X})- \pi\|_{L^2}^2$. Given \Cref{lemma:histogramutility} and the fact that we work on probability distributions with a compact support, upper-bounding this quantity is straightforward.

The lower-bound for the integrated risk is given by :

\begin{theorem}[Integrated lower-bound]
\label{th:lidp}
There exists $C_L>0$, a positive constant depending only on $L$ such that, there exist $n_0(L) \in \N$ and $c_0(L) > 0$ such that for any $n \geq n_0$, and any $\alpha \geq c_0/n$
%such that $n \alpha \geq c_0$,
\begin{equation*}
    \inf_{\hat{\pi} \text{ s.t. } \mathcal{C}} \sup_{\pi \in \Theta^{\text{Lip}}_L}
    \E_{\vect{X} \sim \Prob_{\pi}^{\otimes n}, \hat{\pi}}\|\hat{\pi}(\vect{X})- \pi\|_{L^2}^2 \geq C_L^{-1} \max \left\{ n^{-2/3}, (n \alpha)^{-1} \right\}
%    \;.
\end{equation*}
where $\alpha=\epsilon$
%Here, $\alpha$ should be interpreted s $\epsilon$ 
when $\mathcal{C}$ is the $\epsilon$-DP condition, and 
%as 
$\alpha=\sqrt{\rho}$ when $\mathcal{C}$ is the $\rho$-zCDP condition.
%and, there exist $d_0(L)$ such that, for any $n \geq n_0$, and any $\rho$ such that $n \sqrt{\rho} \geq c_0$ and $\frac{n}{\rho} \geq d_0$,
%\begin{equation*}
%    \inf_{\hat{\pi} \text{ s.t. } \text{$\rho$-zCDP}} \sup_{\pi \in \Theta^{\text{Lip}}_L}
%    \E_{\vect{X} \sim \Prob_{\pi}^{\otimes n}, \hat{\pi}}\|\hat{\pi}(\vect{X})- \pi\|_{L^2}^2 \geq C_L^{-1} \max \left\{ n^{-2/3}, (n \sqrt{\rho})^{-4/3} \right\}
%    \;.
%\end{equation*}
\end{theorem}

\begin{proof}[Proof idea]
If we were to use the same packing %(see \Cref{fig:onetriangle}) 
as in the proof of \Cref{th:lpdp}, the lower-bounds would not be good. Indeed, moving from the pointwise difference to the $L^2$ norm significantly diminishes the distances in the packing. Instead, we will use %\AG{donner une ref où une idée semblable est présente} 
the same idea of deviating from a constant function by triangles, except that we authorize more than one deviation. More specifically, we consider a packing consisting of densities $f_{\omega}$'s where the $\omega$'s are a well-chosen family of $\{ 0, 1 \}^m$ ($m$ is fixed in the proof) \citep{van2000asymptotic}. Then, for a given  $\omega \in \{ 0, 1 \}^m$, $f_{\omega}$ has a triangle centered on $\frac{i}{m+1}$ iff $w_i \neq 0$. %Such a density is represented in \Cref{fig:multipletriangles}.

We then leverage Fano-type inequalities, and we use Assouad's method in order to find the announced lower-bounds.
\end{proof}
The full proof is in \Cref{proofofth:lidp}.

Since the lower-bounds of \Cref{th:lidp} match the upper-bounds of \Cref{proposition:histdp}, we conclude that the corresponding estimators are optimal in terms of minimax rate of convergence.
%\paragraph{New discussion about optimality}
%For the $L^2$-norm, under $\epsilon$-DP, this result leads to the same conclusions as with the pointwise risk (or in infinite norm). Under $\rho$-zCDP on the other hand, we can see that the lower-bound on \Cref{th:lidp} does not quite match the upper bound of \Cref{proposition:histdp}. Nevertheless, \Cref{th:lidp} still says that in low privacy regimes, a degradation of the rate of estimation is to be expected.

\section{Projection Estimators and Periodic Sobolev Densities}
\label{sec:projection}

The Lipschitz densities considered in \Cref{sec:histogram} are general enough to be applicable in many problems. However, this level of generality becomes a curse in terms of rate of estimation. Indeed, as we have seen, the optimal rate of estimation is $\max \p{n^{-2/3}, (n \epsilon)^{-1}}$. To put it into perspective, for many parametric estimation procedures, the optimal rate of convergence usually scales as $\max \p{n^{-1}, (n \epsilon)^{-2}}$ \citep{acharya2021differentially}. This section studies the estimation of smoother distributions, for different smoothness levels, at the cost of generality. In particular, it establishes that the smoother the distribution class is, the closer the private rate of estimation is to $\max \p{n^{-1}, (n \epsilon)^{-2}}$. In other words, it means that the more regular the density is supposed to be, the closer we get to the difficulty of parametric estimation.


When the density of interest $\pi$ is in $L^2([0, 1])$, it is possible to approximate it by projections. Indeed, $L^2([0, 1])$ being a separable Hilbert space, there exists a countable orthonormal family $\p{\phi_i}_{i \in \N \setminus \{0\}}$ that is a Hilbert basis. In particular, if $\theta_i \eqdef \int_{[0, 1]} \pi\;\phi_i $ then
\begin{equation*}
    \sum_{i=1}^N \theta_i \phi_i \stackrel{L^2}{\conv{N \rightarrow +\infty}} \pi \;.
\end{equation*}
Let $N$ be a positive integer, $Z_1, \dots, Z_{N}$ be independent and identically distributed random variables with the same distribution as a centered random variable $Z$ having a finite variance. Given a dataset $\vect{X} = (X_1, \dots, X_n)$, that is also independent of $Z_1, \dots, Z_{N}$, the (randomized) projection estimator is defined as 
\begin{equation}
\label{def:projestimator}
    \hat{\pi}^{\text{proj}} (\vect{X}) = \sum_{i=1}^N \p{\hat{\theta}_i + \frac{1}{n} Z_i} \phi_i \quad \text{where} \quad \hat{\theta}_i \eqdef \frac{1}{n} \sum_{j=1}^n \phi_i(X_j)  \;.
\end{equation}%\mRG{Faire ressortir qq part que $Z_i$ sont aussi indépendantes du dataset, c'est utilisé dans la preuve}
The truncation order $N$ and the random variable $Z$ are tuned later to obtain the desired levels of privacy and utility.
$L^2([0, 1])$ has many well known Hilbert bases, hence multiple choices for the family $\p{\phi_i}_{i \in \N \setminus \{0\}}$. For instance, orthogonal polynomials, wavelets, or the Fourier basis, are often great choices for projection estimators. Because of the privacy constraint however, it is better to consider a \emph{uniformly bounded} Hilbert basis \citep{wasserman2010statistical}, which is typically not the case with a polynomial or wavelet basis. From now on, this work will focus on the following Fourier basis : 
\begin{equation*}
    \begin{aligned}
        \phi_1(x) &= 1 \\
        \phi_{2k}(x) &= \sqrt{2} \sin \p{2 \pi k x} \quad k \geq 1 \\
        \phi_{2k + 1}(x) &= \sqrt{2} \cos \p{2 \pi k x} \quad k \geq 1 \;.
    \end{aligned}
\end{equation*}

\subsection{General utility of projection estimators}

By the Parseval formula, the truncation resulting of approximating the density $\pi$ on a finite family of $N$ orthonormal functions induces a bias term that accounts for $\sum_{i \geq N+1} \theta_i^2$ in the mean square error. Characterizing the utility of $\hat{\pi}^{\text{proj}} $  requires controlling this term, and this is usually done by imposing that $\pi$ is in a Sobolev space. We recall the definition given in \citet{tsybakov2003introduction}: given $\beta \in \N \setminus \{0\}$ and $L>0$, the class $\Theta^{\text{Sob}}_{L, \beta}$ of Sobolev densities of parameters $\beta$ and $L$ is defined as 
\begin{equation*}
\Theta^{\text{Sob}}_{L, \beta} \eqdef \left\{ \pi \in \mathcal{C}^\beta([0, 1], \R_+) \left| 
\begin{cases}
  \pi^{(\beta - 1)} \text{ is absolutely continuous}\;, \\
  \int_{[0, 1]} \p{\pi^{(\beta)}}^2 \leq L^2\;, \\
  \int_{[0, 1]} \pi = 1 \;.  
\end{cases}  
 \right.\right\} \;.
\end{equation*}
For a function $f$, we used the notation $f^{(\beta)}$ to refer to its derivative of order $\beta$.
In addition, the class $\Theta^{\text{PSob}}_{L, \beta}$ of periodic Sobolev densities of parameters $\beta$ and $L$ is defined as 
\begin{equation}
\label{eq:defpsob}
\Theta^{\text{PSob}}_{L, \beta} \eqdef \left\{ \pi \in \Theta^{\text{Sob}}_{L, \beta} \left| 
\forall j \in \{0, \dots, \beta - 1 \}, \pi^{(j)}(0) =  \pi^{(j)}(1)
 \right.\right\} \;.
\end{equation}
Finally, we recall the following general-purpose lemma \citep{tsybakov2003introduction} that allows controlling the truncation bias : 
\begin{fact}[{Ellipsoid reformulation \citep{tsybakov2003introduction}}]
\label[fact]{lemma:sobolevellipsoid}
    A non-negative function $\pi$ with integral $1$ belongs to $\Theta^{\text{PSob}}_{L, \beta}$ if and only if $\displaystyle{\sum_{i=1}^{\infty} a_i^{2 \beta} \theta_i^2 \leq \frac{L^2}{\pi^{2 \beta}}}$, where $a_j \eqdef j$ if $j$ is even and $a_j \eqdef j-1$ if $j$ is odd.
\end{fact}

In this class, one can characterize the utility of projection estimators with the following lemma:
\begin{lemma}[General utility of \eqref{def:projestimator}]
\label[lemma]{lemma:projectionutility}
There is a constant $C_{L, \beta} > 0$, depending only on $L,\beta$, such that
\begin{equation*}
    \sup_{\pi \in \Theta^{\text{PSob}}_{L, \beta}}
    \E_{\vect{X} \sim \Prob_{\pi}^{\otimes n}, \hat{\pi}^{\text{proj}}} \| \hat{\pi}^{\text{proj}}(\vect{X}) - \pi\|_{L^2}^2
    \leq C_{L, \beta} \p{\frac{1}{N^{2 \beta}} + \frac{N}{n} + \frac{N \V(Z)}{n^2}} \;.
\end{equation*}
\end{lemma}

The proof can be found in \Cref{proofoflemma:projectionutility}

\subsection{Privacy and bias tuning}

The estimator $\hat{\pi}^{\text{proj}}(\vect{X})$ is a function of the sums $\p{\sum_{j=1}^n \phi_1(X_j), \dots, \sum_{j=1}^n \phi_N(X_j)}$.  In particular, it is possible to use Laplace and Gaussian mechanisms on this function in order to obtain privacy. Since the functions $|\phi_i|$ are bounded by $\sqrt{2}$ for any $i$, the $l_1$ sensitivity of this function is $2 \sqrt{2} N$ and its $l_2$ sensitivity %\mRG{1) $l_1$ and $l_2$? Cohérence des notations qui étaient introduites comme $l_1$ et $L^2$ en section 2?} 
is $2 \sqrt{2} \sqrt{N}$. Applying the Laplace and the Gaussian mechanism and tuning $N$ %m 
to optimize the utility of \Cref{lemma:projectionutility} gives the following result:

\begin{theorem}[Privacy and utility of \eqref{def:projestimator}] 
\label[theorem]{proposition:projdp}
Given any $\epsilon > 0$ and truncation order $N$,
using $\hat{\pi}^{\text{proj}}$  with  $Z = \frac{2 N \sqrt{2}}{\epsilon} \mathcal{L}(1)$, where $\mathcal{L}(1)$ refers to a random variable following a Laplace distribution of parameter $1$, leads to an $\epsilon$-DP procedure. Moreover, there exists $C_{L, \beta} > 0$, a positive constant that only depends on $L$ and $\beta$, such that if $N$ is on the order of
$\min\p{n^{\frac{1}{2 \beta + 1}}, \p{n \epsilon}^{\frac{1}{\beta + 3/2}}}$,
\begin{equation*}
    \sup_{\pi \in \Theta^{\text{PSob}}_{L, \beta}}
    \E_{\vect{X} \sim \Prob_{\pi}^{\otimes n}, \hat{\pi}^{\text{proj}}} \| \hat{\pi}^{\text{proj}}(\vect{X}) - \pi \|_{L^2}^2
    \leq C_{L, \beta}  \max \left\{ n^{-\frac{2 \beta}{2 \beta + 1}}, (n \epsilon)^{-\frac{2 \beta}{ \beta + 3/2}} \right\} \;.
\end{equation*}
Furthermore, given any $\rho > 0$, and truncation order $N$,
using $\hat{\pi}^{\text{proj}}$ with $Z = \frac{2 \sqrt{N}}{\sqrt{\rho}}\mathcal{N}(0, 1)$, where $\mathcal{N}(0, 1)$ refers to a random variable following a centered Gaussian distribution of variance $1$, leads to a $\rho$-zCDP procedure. Moreover, there exists $C_{L, \beta} > 0$, a positive constant that only depends on $L$ and $\beta$, such that, if $N$ is of the order of 
$\min\p{n^{\frac{1}{2 \beta + 1}}, \p{n \sqrt{\rho}}^{\frac{1}{\beta + 1}}}$
\begin{equation*}
    \sup_{\pi \in \Theta^{\text{PSob}}_{L, \beta}}
    \E_{\vect{X} \sim \Prob_{\pi}^{\otimes n}, \hat{\pi}^{\text{proj}}} \| \hat{\pi}^{\text{proj}}(\vect{X}) - \pi \|_{L^2}^2
    \leq C_{L, \beta}  \max \left\{ n^{-\frac{2 \beta}{2 \beta + 1}}, (n \sqrt{\rho})^{-\frac{2 \beta}{ \beta + 1}} \right\} \;.
\end{equation*}
\end{theorem}
We now discuss these guarantees depending on the considered privacy regime.
\paragraph{Low privacy regimes.} 

According to \Cref{proposition:projdp}, when the privacy-tuning parameters are not too small (i.e. when the estimation is not too private), the usual rate of convergence $n^{-\frac{2 \beta}{2 \beta + 1}}$ is not degraded. In particular, for constant $\epsilon$ or $\rho$, this recovers the results of \citet{wasserman2010statistical}.

\paragraph{High privacy regimes.}

Furthermore, \Cref{proposition:projdp} tells that in high privacy regimes ($\epsilon \ll n^{- \frac{\beta - 1/2}{2 \beta + 1}}$ or $\rho \ll n^{- \frac{2\beta +2}{2 \beta + 1}}$), the provable guarantees of the projection estimator are degraded compared to the usual rate of convergence. Is this degradation constitutive of the estimation problem, or is it due to a suboptimal upper-bound? \Cref{sec:lbsob} shows that this excess of risk is in fact almost optimal.

\subsection{Lower-bounds}

\label{sec:lbsob}

As with the integrated risk on Lipschitz distributions, obtaining lower-bounds for the class of periodic Sobolev densities is done by considering a packing with many elements. The idea of the packing is globally the same as for histograms, except that the uniform density is perturbed with a general $C^\infty$ kernel with compact support instead of simple triangles. In the end, we obtain the following result:

\begin{theorem}[Integrated lower-bound]
\label{th:sidp}
Given $L,\beta>0$ there exists constants
%there exists a constant 
$C_{L, \beta}>0$,
%depending only on $L$ and $\beta$ such that, there exist 
$n_0(L, \beta) \in \N$, and $c_0(L, \beta) > 0$, such that for any $n \geq n_0$, and any $\alpha \geq c_0/n$ %such that $n \alpha \geq c_0$,
\begin{equation*}
    \inf_{\hat{\pi} \text{ s.t. } \mathcal{C}} \sup_{\pi \in \Theta^{\text{PSob}}_{L, \beta}}
    \E_{\vect{X} \sim \Prob_{\pi}^{\otimes n}, \hat{\pi}}\|\hat{\pi}(\vect{X})- \pi\|_{L^2}^2 \geq  C_{L, \beta}^{-1} \max \left\{ n^{-\frac{2 \beta}{2 \beta + 1}}, (n \alpha)^{-\frac{2 \beta}{\beta + 1}} \right\}
 %   \;.
\end{equation*}
where $\alpha=\epsilon$
%ere, $\alpha$ should be interpreted as $\epsilon$ 
when $\mathcal{C}$ is the $\epsilon$-DP condition, and 
%as 
$\alpha=\sqrt{\rho}$ when $\mathcal{C}$ is the $\rho$-zCDP condition.
%and, there exist $d_0(L, \beta)$ such that, for any $n \geq n_0$, and any $\rho$ such that $n \sqrt{\rho} \geq c_0$ and $\frac{n}{\p{n \sqrt{\rho}}^{\frac{2 \beta}{2 \beta + 1}}} \geq d_0$,
%\begin{equation*}
%    \inf_{\hat{\pi} \text{ s.t. } \text{$\rho$-zCDP}} \sup_{\pi \in \Theta^{\text{PSob}}_{L, \beta}}
%    \E_{\vect{X} \sim \Prob_{\pi}^{\otimes n}, \hat{\pi}}\| \hat{\pi}(\vect{X})- \pi \|_{L^2}^2 \geq      C_{L, \beta}^{-1} \max \left\{ n^{-\frac{2 \beta}{2 \beta + 1}}, (n \sqrt{\rho})^{-\frac{2 \beta}{\beta + 1/2}} \right\}
%    \;.
%\end{equation*}
\end{theorem}

\begin{proof}[Proof idea]
As with the proof of \Cref{th:lidp}, this lower-bound is based on the construction of a packing of densities $f_{\omega}$'s where the $\omega$'s are a well-chosen family of $\{ 0, 1 \}^m$ ($m$ is fixed in the proof). Then, for a given  $\omega \in \{ 0, 1 \}^m$, $f_{\omega}$ deviates from a constant function around $\frac{i}{m+1}$ if, and only if, %iff
$w_i \neq 0$. 
Contrary to the proof of \Cref{th:lidp} however, the deviation cannot be by a triangle : Indeed, such a function wouldn't even be differentiable. Instead, we use a deviation by a $C^\infty$ kernel with compact support. Even if the complete details are given in the full proof.% \Cref{fig:multiplekernels} gives a general illustration of the packing. 

Again, Fano-type inequalities (for the $\epsilon$-DP case), and Assouad's lemma (for the $\rho$-zCDP case) are used to conclude.
\end{proof}

The full proof can be found in \Cref{proofofth:sidp}.
In comparison with the upper-bounds of \Cref{proposition:projdp}, for $\epsilon$-DP the lower-bound \emph{almost} matches the guarantees of the projection estimator. In particular, the excess of risk in the high privacy regime is close to being optimal. \Cref{sec:relaxation} explains how to bridge the gap even more, at the cost of relaxation. 

Under $\rho$-zCDP,  %on the other hand, 
the lower-bounds and upper-bounds actually match. We conclude that projection estimators with $\rho$-zCDP obtain minimax-optimal rates of convergence.

\subsection{Near minimax optimality via relaxation}
\label{sec:relaxation}

An hypothesis that we can make on the sub-optimality of the projection estimator against $\epsilon$-DP mechanisms is that the $l_1$%\mRG{$l_1$, $l1$, $L^1$? Idem $l2$ \ldots} 
sensitivity of the estimation of $N$ Fourier coefficients scales as $N$ whereas its $l_2$ sensitivity scales as $\sqrt{N}$. Traditionally, the Gaussian mechanism \citep{dwork2006our,dwork2006calibrating} has allowed to use the $l_2$ sensitivity instead of the $l_1$ one at the cost of introducing a relaxation term $\delta$ in the privacy guarantees, leading to $(\epsilon, \delta)$-DP.

%\citet{bun2016concentrated} states that if a mechanism $\mech{M}$ is $\rho$-zCDP, then it is $\p{\rho + 2 \sqrt{ \rho \ln(1/\delta)}, \delta}$-DP for any $\delta > 0$. Applying this result
%\mRG{Donner une preuve explicite meme courte en annexe? \CL{C'est vraiment une application directe de Bun et Steinke}} 
%and \Cref{proposition:projdp} with $\delta$ of the form $\frac{1}{n^\gamma}$ for a positive exponent $\gamma$ gives the following result :  

\begin{corollary}[Privacy and utility of \eqref{def:projestimator} with relaxation] 
\label{proposition:relaxation}
%\mRG{Hypotheses a malaxer un peu pour clarifier}
Consider $\gamma>0$, $n$ and integer, and $0<\epsilon \leq 8 \ln n^\gamma$.
%For any $\epsilon > 0$ and $\gamma > 0$, d
Defining $\Tilde{\rho} \eqdef \frac{1}{16} \frac{\epsilon^2}{\ln\p{n^\gamma}}$
and using $\hat{\pi}^{\text{proj}}$ with $Z = \frac{2 \sqrt{N}}{\sqrt{\Tilde{\rho}}}\mathcal{N}(0, 1)$, 
where $\mathcal{N}(0, 1)$ refers to a random variable following a centered Gaussian distribution of variance $1$, 
leads to an $\p{\epsilon, \frac{1}{n^\gamma}}$-DP procedure. % if $\epsilon \leq 8 \ln\p{n^\gamma}$. Furthermore, 
%$N$ is of the order of  $\min\p{n^{\frac{1}{2 \beta + 1}}, \p{n \sqrt{\Tilde{\rho}}}^{\frac{1}{\beta + 1}}}$, 
there exists $C_{L, \beta} > 0$, a positive constant that only depends on $L$ and $\beta$, such that if $N$ is of the order of  
%$\min\p{n^{\frac{1}{2 \beta + 1}}, \p{\frac{n}{4} \frac{\epsilon}{\sqrt{\ln\p{n^{\gamma}}}}}^{\frac{1}{\beta + 1}}}$ 
%
%vu qu'il s'agit d'ordres on peut écrire
$\min\p{n^{\frac{1}{2 \beta + 1}}, \p{\frac{n}{\sqrt{\ln n}} \cdot \frac{\epsilon}{\sqrt{\gamma}}}^{\frac{1}{\beta+1}}}$
then
\begin{equation*}
\sup_{\pi \in \Theta^{\text{PSob}}_{L, \beta}}
    \E_{\vect{X} \sim \Prob_{\pi}^{\otimes n}, \hat{\pi}^{\text{proj}}} \| \hat{\pi}^{\text{proj}}(\vect{X}) - \pi \|_{L^2}^2
    \leq C_{L, \beta}  \max \left\{ n^{-\frac{2 \beta}{2 \beta + 1}}, P_{\beta, \gamma}(\ln(n)) (n \epsilon)^{-\frac{2 \beta}{ \beta + 1}} \right\} \;,
\end{equation*}
where $P_{\beta, \gamma}$ is a polynomial expression depending on $\beta$ and $\gamma$.
\end{corollary}
\begin{proof}
    %Indeed, when 
    Since $\epsilon \leq 8 \ln\p{n^\gamma}$, we have $\Tilde{\rho} \leq 2 \sqrt{ \Tilde{\rho} \ln(n^\gamma)}$. By \Cref{proposition:projdp} the mechanism is $\tilde{\rho}$-zCDP, and satisfies the claimed upper bounds for $N$ on the considered order. By \citet{bun2016concentrated} (that states that if a mechanism $\mech{M}$ is $\rho$-zCDP, then it is $\p{\rho + 2 \sqrt{ \rho \ln(1/\delta)}, \delta}$-DP for any $\delta > 0$) it is thus
    %The resulting mechanism is thus
    $\p{4 \sqrt{ \Tilde{\rho} \ln(n^{\gamma})},  \frac{1}{n^\gamma}}$-DP.
\end{proof}
%\RG{Pour moi on a $\tilde{\rho} \ln 1/\delta = \frac{\epsilon^{2}}{16} (\ln n^{\gamma})^{2}$ donc Bun \& Steinke donne ici un mécanisme $(\tilde{\epsilon},\delta)$-DP avec $\tilde{\epsilon} \coloneqq \tilde{\rho}+\frac{\epsilon}{2} \ln n^{\gamma} > \epsilon$ dès que $n^{\gamma} \geq e^{2}$. Pour obtenir $(\epsilon,\delta)$-DP ne peut-on pas juste résoudre $\tilde{\rho}+2\sqrt{\tilde{\rho} \ln1/\delta} = \epsilon$, cad $(\tilde{\rho}-\epsilon)^{2} = 4 \tilde{\rho} \ln1/\delta$, et conclure avec un bete trinome. Si lexpression est trop laide, prendre $\rho' < \tilde{\rho}$ avec une expression plus simple doit suffire?  \CL{Il y avait effectivement une erreur de calcul, je me suis chargé d'écrire une preuve.}}
In order to understand the implications of this result, one must understand the role of $\delta$ in $(\epsilon, \delta)$-differential privacy. It is usually interpreted as the probability of the procedure not respecting the $\epsilon$-DP condition \citep{dwork2014algorithmic}. Hence, with probability $\delta$, the result is not guaranteed to be private. A general rule of thumb for choosing $\delta$ is to take it much smaller than $1/n$ so that each individual of the database only has a small chance of seeing its data leak \citep{dwork2014algorithmic}. Choosing $\delta = 1/n^\gamma$ for $\gamma > 1$ is hence considered a good choice for $\delta$.

With this relaxation, the upper-bound of \Cref{proposition:relaxation} matches the lower-bound of \Cref{th:sidp} for $\epsilon$-DP up to polylog factors.

\section{Conclusion}\label{sec:ccl}

As we have seen throughout this article, under global privacy, one can usually distinguish two estimation regimes. In the \emph{low} privacy regime, on the one hand, the estimation rate is not degraded compared to its non-private counterpart. This notably covers the early observation of \citet{wasserman2010statistical} for constant privacy budget. 
In the \emph{high} privacy regime on the other hand, a provable degradation is unavoidable,  %to be expected. 
and we extended the study of such regimes beyond the cases covered in \citet{barber2014privacy}.

Besides examples in which the estimation is sharp in both regimes, we also presented some example in which there are \emph{small} gaps between the proved upper-bounds and lower-bounds. These gaps are nevertheless very small, especially for high degrees of smoothness, and they can be bridged up to logarithmic factors with a \emph{reasonable} and quite standard relaxation.

%Nevertheless, this article demonstrated that there is still some room for improvement on the topic of density estimation under \emph{global} differential privacy, even without considering advanced properties such as adaptivity. We hope that it will kick-start a regain of interest for the problem.

%\subsection*{Acknowledgement}
%Aurélien Garivier acknowledges the support of the Project IDEXLYON of the University of Lyon, in the framework of the Programme 
%Investissements d'Avenir (ANR-16-IDEX-0005), and Chaire SeqALO (ANR-20-CHIA-0020-01).
%This project was supported in part by the AllegroAssai ANR project ANR-19-CHIA-0009.
%Additionally, we thank the anonymous reviewers for their precious inputs and suggestions.

\subsection*{Acknowledgement}
Aurélien Garivier acknowledges the support of the Project IDEXLYON of the University of Lyon, in the framework of the Programme 
Investissements d'Avenir (ANR-16-IDEX-0005), and Chaire SeqALO (ANR-20-CHIA-0020-01).
This project was supported in part by the AllegroAssai ANR project ANR-19-CHIA-0009.
Additionally, we thank the anonymous reviewers for their precious inputs and suggestions.

%\newpage
\bibliography{biblio}
\bibliographystyle{tmlr}

\newpage
\appendix

\section{Useful results from the litterature}
\label{section:usefullemmas}

\begin{fact}[{Neyman-Pearson \& Le Cam's lemma \citep[Lemma 5.3]{rigollet2015high}}]
\label[fact]{fact:lecamslemma}
Let $\Prob_1,\\ \Prob_2$ be two probability distributions on a measure space $\set{E}$, then
\begin{equation}
     \begin{aligned}
        \inf_{\Psi : \set{E}  \rightarrow \{1, 2\}} \max_{i \in \{1, 2\}} \Prob_{\mathbf{X} \sim \Prob_i} \p{\Psi \p{\vect{X}} \neq i}  
          &\geq \frac{1}{2} \inf_{\Psi : \set{E}  \rightarrow \{1, 2\}} \sum_{i=1}^2 \Prob_{\mathbf{X} \sim \Prob_i} \p{\Psi \p{\vect{X}} \neq i} \\
          &= \frac{1}{2} \p{1 - \tv{\Prob_1}{\Prob_2}} \;.
     \end{aligned}
     \label{eq:ClassicalLeCam}
 \end{equation}
\end{fact}

\begin{fact}[{Fano's lemma \citep[Theorem 3.1]{giraud2021introduction}}]
\label[fact]{fact:fanoslemma}
 Let $\p{\Prob_i}_{i \in \intslice{1}{N}}$ be a family of probability distributions on a measure space $\set{E}$. For any probability distribution $\mathbb{Q}$ on $\set{E}$ such that $\Prob_i \ll \mathbb{Q}$ for all $i$, and for any test function $\Psi : \set{X}^n  \rightarrow \intslice{1}{N}$, 
 \begin{equation}
     \begin{aligned}
         \max_{i \in \intslice{1}{N}} \Prob_{\mathbf{X} \sim \Prob_i} \p{\Psi \p{\vect{X}} \neq i}  
          &\geq \frac{1}{N} \sum_{i=1}^N \Prob_{\mathbf{X} \sim \Prob_i} \p{\Psi \p{\vect{X}} \neq i} \\
          &\geq 1 - \frac{1 + \frac{1}{N} \sum_{i=1}^N \kl{\Prob_i}{\mathbb{Q}}}{\ln(N)} \;.
     \end{aligned}
     \label{eq:ClassicalFano}
 \end{equation}
 Often $\mathbb{Q}$ is set to $\frac{1}{N} \sum_{i=1}^N \Prob_i$.
\end{fact}

\begin{fact}[{Le Cam's lemma for differential privacy \cite[Theorem 1]{lalanne2022statistical}}]
\label[fact]{fact:lecamslemmaprivate}
    If a randomized mechanism $\mech{M}$ satisfies $(\epsilon, \delta)$-DP, then for any test function $\Psi : \codom{\mech{M}} \rightarrow \{1, 2\}$ and any probability distributions $\Prob_1$ and $\Prob_2$ on $\set{X}$  we have
\begin{equation*}
\begin{aligned}
    \max_{i \in \{1, 2\}} \Prob_{\mathbf{X} \sim \Prob_i^{\otimes n}, \mech{M}} &\p{\Psi \p{\mech{M}\p{\vect{X}}} \neq i} \\
    \geq 
    \frac{1}{2} 
    & \p{\p{1 - \left( 1 - e^{-\epsilon} \right) \tv{\Prob_{1}}{\Prob_{2}}  }^n -  2 n e^{-\epsilon} \delta \tv{\Prob_{1}}{\Prob_{2}}} \;.
\end{aligned}
\end{equation*}
\end{fact}


\begin{fact}{Le Cam's lemma for concentrated differential privacy \cite[Theorem 2]{lalanne2022statistical}}]
\label[fact]{fact:lecamslemmaconcentrated}
If a randomized mechanism $\mech{M}$ satisfies $\rho$-zCDP, then for any test function $\Psi : \codom{\mech{M}}  \rightarrow \intslice{1}{N}$ and any probability distributions $\Prob_1$ and $\Prob_2$ on $\set{X}$,
\begin{equation*}
\begin{aligned}
    \max_{i \in \{1, 2\}} \Prob_{\mathbf{X} \sim \Prob_i^{\otimes n}, \mech{M}} \p{\Psi \p{\mech{M}\p{\vect{X}}} \neq i} 
    \geq 
    \frac{1}{2} 
    & \p{1 - n \sqrt{\rho / 2}  \tv{\Prob_1}{\Prob_2}} \;.
\end{aligned}
\end{equation*}
\end{fact}

\begin{fact}[{Fano's lemma for differential privacy \cite[Theorem 3]{lalanne2022statistical}}]
\label[fact]{fact:fanoslemmaprivate}
If a randomized mechanism $\mech{M}$ satisfies $\epsilon$-DP, then for any test function $\Psi : \codom{\mech{M}}  \rightarrow \intslice{1}{N}$, any family of probability distributions $\p{\Prob_i}_{i \in \intslice{1}{N}}$ on $\set{X}$,
\begin{equation*}
     \begin{aligned}
         \max_{i \in \intslice{1}{N}} \Prob_{\mathbf{X} \sim \Prob_i^{\otimes n}, \mech{M}} \p{\Psi \p{\mech{M}(\vect{X})} \neq i}  
          \geq  1 - \frac{1 + \frac{n \epsilon}{N^2} \sum_{i, j} \frac{2 \tv{\Prob_i}{\Prob_j}}{1 + \tv{\Prob_i}{\Prob_j}}}{\ln(N)}    \;. \\
     \end{aligned}
 \end{equation*}
\end{fact}

\begin{fact}[{Fano's lemma for differential privacy \cite[Theorem 4]{lalanne2022statistical}}]
\label[fact]{fact:fanoslemmaconcentrated}
If a randomized mechanism $\mech{M}$ satisfies $\epsilon$-DP, then for any test function $\Psi : \codom{\mech{M}}  \rightarrow \intslice{1}{N}$, any family of probability distributions $\p{\Prob_i}_{i \in \intslice{1}{N}}$ on $\set{X}$,
 \begin{equation*}
     \begin{aligned}
         \max_{i \in \intslice{1}{N}} \Prob_{\mathbf{X} \sim \Prob_i^{\otimes n}, \mech{M}} \p{\Psi \p{\mech{M}(\vect{X})} \neq i} 
          &\geq 1 - \frac{1 + \frac{n^2 \rho}{N^2} \sum_{i, j}\frac{1}{n} \frac{2  \tv{\Prob_i}{\Prob_j}}{1 + \tv{\Prob_i}{\Prob_j}} + \p{\frac{2  \tv{\Prob_i}{\Prob_j}}{1 + \tv{\Prob_i}{\Prob_j}}}^2}{\ln(N)}    \;.
     \end{aligned}
 \end{equation*}
\end{fact}

%\section{Figures}
%\label{appendix:figures}

%\begin{figure}[h]
%    \centering
%    \includesvg[scale=0.5]{fig1.svg}
%   \caption{Packing for \Cref{th:lpdp}}
%    \label{fig:onetriangle}
%\end{figure}

%\begin{figure}[h]
%    \centering
%    \includesvg[scale=0.5]{fig2.svg}
%    \caption{Packing for \Cref{th:lidp}}
%    \label{fig:multipletriangles}
%\end{figure}

%\begin{figure}[h]
%\centering
%\includesvg[scale=0.5]{fig3.svg}
%\caption{Packing for \Cref{th:sidp}}
%\label{fig:multiplekernels}
%\end{figure}

\section{Proof of \Cref{lemma:histogramutility}}
\label{proofoflemma:histogramutility}

Let $\pi \in \Theta^{\text{Lip}}_L$, $x_0 \in [0, 1]$. The classical bias-variance decomposition gives that%\mRG{Ici et ailleurs, utiliser $\{[()]\}$ pour une meilleure lisibilité des parenthèsages imbriqués?}
\begin{equation*}
\begin{aligned}
    \E&\p{ \p{\hat{\pi}^{\text{hist}}(\vect{X})(x_0) - \pi(x_0)}^2}
    = \p{\E \p{\hat{\pi}^{\text{hist}}(\vect{X})(x_0)} - \pi(x_0)}^2
    + \V \p{\hat{\pi}^{\text{hist}}(\vect{X})(x_0)} \;.
\end{aligned}
\end{equation*}
For any $x \in [0, 1]$, we note $\text{bin}(x)$ the bin of the histogram in which $x$ falls into. Notice that, for any $x_0 \in [0,1]$ and any integer $i$, the random variable $\Ind_{ \text{bin}(x_0)}(X_i)$ follows a Bernoulli distribution of probability of success $\int_{\text{bin}(x_0)} \pi$. Let us first study the bias, using the definition~\eqref{eq:hist} of $\hat{\pi}^{\text{hist}}$
\begin{equation*}
    \begin{aligned}
        &\left| \E\p{\hat{\pi}^{\text{hist}}(\vect{X})(x_0)} - \pi(x_0) \right|
        =
        \left| \frac{1}{nh} \sum_{i = 1}^n \E\p{\Ind_{ \text{bin}(x_0)}(X_i)}   - \pi(x_0) \right|
        =
        \left| \frac{n\int_{\text{bin}(x_0)} \pi(x)dx}{nh} - \pi(x_0) \right| \\
        &=
        \frac{1}{h} \left| \int_{\text{bin}(x_0)} (\pi(x) - \pi(x_0)) dx \right| 
        \leq 
        \frac{1}{h} \int_{\text{bin}(x_0)} \left|  \pi(x) - \pi(x_0) \right| dx
        \leq 
        \frac{L}{h} \int_{\text{bin}(x_0)} \left|  x - x_0 \right| dx
        \leq 
        \frac{L h}{2} \;.
    \end{aligned}
\end{equation*}
Let us now look at the variance. By independence of $X_i$'s and $Z_j$'s, 
\begin{equation*}
    \begin{aligned}
        \V\p{\hat{\pi}^{\text{hist}}(\vect{X})(x_0)} 
        &=
        \frac{1}{n^2h^2} \p{\sum_{i = 1}^n \V\p{\Ind_{ \text{bin}(x_0)}(X_i)} + \V \p{Z_{\text{bin}(x_0)}} } \\
        &=
        \frac{1}{n^2h^2} \p{n \p{\int_{\text{bin}(x_0)} \pi }\p{1-\int_{\text{bin}(x_0)} \pi} + \V \p{Z} } \\ 
        &\leq \frac{1}{n h^2} \p{\int_{\text{bin}(x_0)} \pi } + \frac{\V \p{Z}}{n^2h^2}.
    \end{aligned}
\end{equation*}
Since $\pi$ is $L$-Lipschitz on $[0, 1]$ and has to integrate to $1$ (because it is a density), $\pi$ is uniformly bounded from above by $L+1$ on $[0, 1]$. Hence, $\int_{\text{bin}(x_0)} \pi \leq (L+1) h$ and the result follows.

%\RG{Cela vaudrait-il (ou non?) le coup de faire ressortir explicitement le role de L dans la borne, sous la forme $\frac{L^2}{4}h^2 + \frac{L+1}{nh}+\frac{\V\p{Z}}{n^2 h^2}$ ? Car la constant de Lipschitz ne joue pas devant la variance de $Z$. L'expression n'est qd meme pas bcp plus compliquéee et on peut passer à $C_L$ dans la suite \ldots }

\section{Assouad's lemma with concentrated differential privacy.}
\label{assouadcdp}

As the reduction to a testing problem between multiple hypotheses, Assouad's lemma relies on similar ideas, where the packing has to be parametrized by a hypercube. Its advantage over tools like Fano's lemma is that it only makes tests between pairs of hypotheses (instead of all of them at the same time). The cost of this is that the control of the packing is slightly more difficult.

Suppose that the set of distributions of interest $\mathcal{P}$ contains a family of distributions $(\Prob_\omega)_{\omega \in \{0, 1 \}^m}$ for a certain positive integer $m$. If the loss function (taken quadratic for simplicity) can be decomposed as 
\begin{equation}
\label{eq:Assouadcondition}
    \forall \omega, \omega' \in \{0, 1 \}^m, \quad \| f_\omega - f_{\omega '} \|_{L^2}^2
    \geq 
    2 \tau \sum_{i=1}^m \Ind_{\omega_{i} \neq \omega_{i}'} = 2 \tau\ham{\omega}{\omega'}\;,
\end{equation}
where for any $\omega$, $f_\omega$ represents the density of $\Prob_\omega$,
then the minimax risk can be lower-bounded as (the proof is classical and can be found in \citet[Section 5.4]{acharya2021differentially})
\begin{equation}
\label{eq:Assouadprivate}
\begin{aligned}
    \inf_{\hat{\pi} \text{ s.t. } \mathcal{C}} &\sup_{\Prob \in \mathcal{P}}  \E_{\vect{X} \sim \Prob, \hat{\pi}} (\|\hat{\pi}(\vect{X}) - \pi \|_{L^2}^2) \\
    &\geq 
    \frac{\tau}{16} \sum_{i=1}^m \underset{\Psi : \codom{\mech{M}} \rightarrow \{0, 1\}}{\inf_{\mech{M} \text{ s.t. } \mathcal{C}}} \Prob_{\mathbf{X} \sim \Prob_{\omega^{i, 0}}^{\otimes n}, \mech{M}} \p{ \Psi \p{\mech{M}(\vect{X})} \neq 0} + \Prob_{\mathbf{X} \sim \Prob_{\omega^{i, 1}}^{\otimes n}, \mech{M}} \p{ \Psi \p{\mech{M}(\vect{X})} \neq 1}\;.
\end{aligned}
\end{equation}
where $\Prob_{\omega^{i, 0}}$ and $\Prob_{\omega^{i, 1}}$ are the \emph{mixture} distributions
\begin{equation}\label{eq:AssouadMixtures}
    \Prob_{\omega^{i, 0}} \eqdef \frac{1}{2^{m-1}} \sum_{\omega \in \{0, 1 \}^m | \omega_i = 0} \Prob_\omega 
    \quad \text{ and }
    \quad 
    \Prob_{\omega^{i, 0}} \eqdef \frac{1}{2^{m-1}} \sum_{\omega \in \{0, 1 \}^m | \omega_i = 1} \Prob_\omega \;.
\end{equation}
The term 
\begin{equation*}
    \Prob_{\mathbf{X} \sim \Prob_{\omega^{i, 0}}^{\otimes n}, \mech{M}} \p{ \Psi \p{\mech{M}(\vect{X})} \neq 0} + \Prob_{\mathbf{X} \sim \Prob_{\omega^{i, 1}}^{\otimes n}, \mech{M}} \p{ \Psi \p{\mech{M}(\vect{X})} \neq 1}
\end{equation*}
characterizes the \emph{testing} difficulty between $\Prob_{\omega^{i, 0}}$ and $\Prob_{\omega^{i, 1}}$. It can be controlled by Le Cam's lemma, and by its variants when working under privacy (see \citet{acharya2021differentially,lalanne2022statistical} for differential privacy and \citet{lalanne2022statistical} for concentrated differential privacy). Such results are reminded in \Cref{section:usefullemmas}.

\section{Proof of \Cref{th:lpdp}}
\label{proofofth:lpdp}

    Let $x_0 \in (0, 1)$. As explained in the sketch of the proof, we build a packing consisting of two elements, and after controlling quantities such as their KL divergences or their TV distances, we leverage Le Cam-type inequalities in order to obtain lower-bounds.
    \paragraph{Packing construction.}
    We define the functions $f_{L, x_0, h}, \forall h>0$ as 
\begin{equation}
\label{eq:deftrianglefunction}
    \forall x \in [0, 1], \quad f_{L, x_0, h}(x) 
    \eqdef
    \begin{cases}
      1-Lh^2 \text{ if } x \in [0, x_0-h) \cup [x_0+h, 1],\\
      1-Lh^2+Lh + L (x-x_0) \text{ if } x \in [x_0-h, x_0) \;. \\
      1-Lh^2+Lh - L (x-x_0) \text{ if } x \in [x_0, x_0+h) \\
    \end{cases}
\end{equation}
Note that as soon as $h \leq \min \{ x_0, 1-x_0 \}$, $f_{L, x_0, h} \in \Theta^{\text{Lip}}_L$. The case $x_0 \in \{0, 1 \}$ is treated in the exact same fashion, but by considering functions that only contain "half of a spike" centered on $x_0$. Furthermore, let us note $g$ the function that is constant to $1$ on $[0, 1]$ (we have $g \in \Theta^{\text{Lip}}_L$). 

We start by %redefining
recalling the total variation distance between two probability distributions, and we recall some useful alternative expressions that are used in the proofs of this article.
Given $(\set{U}, \set{T})$ a set $\set{U}$ equipped with a $\sigma$-algebra $\set{T}$, and two probability measures $\Prob_1$ and $\Prob_2$ two probability distributions on $\set{U}$, and compatible with $\set{T}$, the total variation distance $\tv{\cdot}{\cdot}$ between $\Prob_1$ and $\Prob_2$ is defined as 
\begin{equation*}
    \tv{\Prob_1}{\Prob_2} \eqdef \sup_{\set{S} \in \set{T}} |\Prob_1(\set{S}) - \Prob_2(\set{S}) | \;.
\end{equation*}
Furthermore, when $\Prob_1, \Prob_2$ are dominated by a common $\sigma$-finite measure $\mu$ on $(\set{U}, \set{T})$, by noting $p_1 \eqdef \frac{d \Prob_1}{ d \mu }$ and $p_2 \eqdef \frac{d \Prob_2}{ d \mu }$, the Radon-Nikodym derivatives of $\Prob_1$ and $\Prob_2$ with respect to $\mu$, the following alternative expressions to the total variation can be useful :
\begin{equation*}
\begin{aligned}
    \tv{\Prob_1}{\Prob_2} &\eqdef \sup_{\set{S} \in \set{T}} |\Prob_1(\set{S}) - \Prob_2(\set{S}) | %\\
    %&
    = \Prob_1(\{ p_1 > p_2 \}) - \Prob_2(\{ p_1 > p_2 \}) \\
    &= \int_{\{ p_1 > p_2 \}} \p{p_1 - p_2} d \mu %\\
    %&
    = \int_{\{ p_2 \geq p_1 \}} \p{p_2 - p_1} d \mu \\
    &= \frac{1}{2} \int_{\set{U}} |p_1 - p_2| d \mu %\\
    %&
    = 1 - \int_{\set{U}} \min(p_1, p_2) d \mu\;.
\end{aligned}
\end{equation*}
These expressions simply come from considering the events $\{ p_1 > p_2 \}$ and $\{ p_2 \geq p_1 \}$ that form a partition of $\set{U}$, and from the relation $|a - b| = a + b - 2 \min(a, b)$ for any real numbers $a$ and $b$.

Jumping back to our original proof, when $f_{L, x_0, h} \in \Theta^{\text{Lip}}_L$, we can compute the total variation between $\Prob_{f_{L, x_0, h}}$ and  $\Prob_{g}$ the distributions of probability with densities $f_{L, x_0, h}$ and $g$ with respect to Lebesgue's measure on $[0, 1]$, 
\begin{equation}
\label{eq:jnbcdshgvdy}
    \tv{\Prob_{f_{L, x_0, h}}}{\Prob_{g}} 
    = 
    1 - \int_{[0, 1]} \min \p{f_{L, x_0, h}, g}
    \stackrel{\text{Constant part}}{\leq }
    1 - \int_{[0, 1]} 1 - Lh^2 dx = Lh^2 \;.
\end{equation}

Another important measure of discrepancy between probability distributions is the so-called Kullback-Leibler (KL) divergence. For two probability distributions $\mathbb{P}$ and $\mathbb{Q}$ such that $\mathbb{P} \ll \mathbb{Q}$ (absolute continuity),  it is defined as
\begin{equation*}
    \kl{\mathbb{P}}{\mathbb{Q}} = \int \ln \p{\frac{d \mathbb{P}}{ d \mathbb{Q}}} d\mathbb{P} \;.
\end{equation*}

Back to our problem, for $h$ in a neighborhood of $0$, we also have the following Taylor expansion on their KL divergence:
\begin{equation}
\label{eq:dsqlkvjdfhvnqgzsera}
\begin{aligned}
    \kl{\Prob_{g}}{\Prob_{f_{L, x_0, h}}} 
    &= 
    \int_{[0, 1]}\ln \p{\frac{g}{f_{L, x_0, h}}} g 
    =
    \ln\p{\frac{1}{1-Lh^2}} (1-2h) + 2 \int_0^h \ln \p{\frac{1}{1-Lh^2 + Lt} } dt \\
    &\leq
    C \p{h^3 + O(h^4)} \;,
\end{aligned}
\end{equation}
where $C$ is a positive constant depending only on $L$ the $O$ only hides constant factors.
Furthermore, $|g(x_0) - f_{L, x_0, h}(x_0)| = L |h^2-h|$ and $\{g, f_{L, x_0, h}\}$ is thus a $\frac{L }{2} |h^2-h|$ packing of $\Theta^{\text{Lip}}_L$ w.r.t the seminorm $f, g \mapsto \|f-g\| \coloneqq |f(x_0) - g(x_0)|$.
\paragraph{Recovering the usual lower-bound}
By the classical minimax reduction as hypothesis testing \Cref{jkhgqskdjhfbgkjqhwgsdkfjhgkqwsjhdgfjkhwsgdb},
\begin{equation}
\label{eq:ghbbvcdjyuhsgdbefr}
    \begin{aligned}
    \inf_{\hat{\pi}  \text{ s.t. } \mathcal{C}} &\sup_{\pi \in \Theta^{\text{Lip}}_L}
    \E_{\vect{X} \sim \Prob_{\pi}^{\otimes n}, \hat{\pi}}\p{\p{\hat{\pi}(\vect{X})(x_0) - \pi(x_0)}^2} \\
    &\geq 
    \frac{L^2 }{4} \p{h^2-h}^2
    \inf_{\hat{\pi}  \text{ s.t. } \mathcal{C}} \inf_{\Psi : \Theta^{\text{Lip}}_L \rightarrow \{0, 1\}}
    \max  \left\{ \Prob_{\vect{X} \sim \Prob_{g}^{\otimes n}, \hat{\pi}}                \p{\Psi(\hat{\pi}(\vect{X})) \neq 0},  \right.\\
        &\quad\quad\quad\quad\quad\quad\quad\quad\quad\quad\quad\quad\quad\quad\quad\quad\quad\quad\quad\quad \left. \Prob_{\vect{X} \sim \Prob_{f_{L, x_0, h}}^{\otimes n}, \hat{\pi}} \p{\Psi(\hat{\pi}(\vect{X})) \neq 1}\right\} \\
    &\stackrel{\text{\Cref{fact:lecamslemma}}}{\geq}
    \frac{L^2}{8}h^2 (1-h)^2 %(h^2-h)^2
    \p{1 - \tv{\Prob_{g}^{\otimes n}}{\Prob_{f_{L, x_0, h}}^{\otimes n}}}  \\
    &\stackrel{\text{Pinsker}}{\geq} \frac{L^2}{8}h^2 (1-h)^2 %(h^2-h)^2
    \p{1 - \sqrt{\kl{\Prob_{g}^{\otimes n}}{\Prob_{f_{L, x_0, h}}^{\otimes n}}/2}} \\
    &\stackrel{\text{Tensorization}}{=} 
    \frac{L^2}{8} h^2 (1-h)^2 %(h^2-h)^2
    \p{1 - \sqrt{n\kl{\Prob_{g}}{\Prob_{f_{L, x_0, h}}}/2}} \\
    &\stackrel{\eqref{eq:dsqlkvjdfhvnqgzsera}}{\geq} 
    \frac{L^2}{8}h^2 (1-h)^2 %(h^2-h)^2
    \p{1 - \sqrt{\frac{n}{2}\p{\frac{h^3 L^2}{3} + O(h^4)}}} \;.
    \end{aligned}
\end{equation}
The second inequality comes from the so-called Le Cam's lemma \cite{rigollet2015high} that lower-bounds the testing difficulty (without further constraints) between two distributions. The next inequality comes from the so-called Pinsker's inequality \cite{tsybakov2003introduction}, that states that for two probability distributions $\mathbb{P}$ and $\mathbb{Q}$, $\tv{\mathbb{P}}{\mathbb{Q}} \leq \sqrt{\kl{\mathbb{P}}{\mathbb{Q}}/2}$. The last inequality is the result of the so-called tensorization property of the KL divergence that states that for two probability distributions $\mathbb{P}$ and $\mathbb{Q}$, and for an integer $n \geq 1$,
$\kl{\mathbb{P}^{\otimes n}}{\mathbb{Q}^{\otimes n}} \leq n \kl{\mathbb{P}}{\mathbb{Q}}$.


When possible (i.e. when $n$ is big enough), setting $h = \p{\frac{1}{4nL^2}}^{1/3}$ leads to, for $n$ big enough (so that $1-h \geq 1/2$ %$h - h^2 \geq h/2$ 
and $|O(h^4)| \leq \frac{h^3L^2}{3}$),
\begin{equation*}
    \inf_{\hat{\pi} \text{ s.t. } \mathcal{C}} \sup_{\pi \in \Theta^{\text{Lip}}_L}
    \E_{\vect{X} \sim \Prob_{\pi}^{\otimes n}, \hat{\pi}}\p{\p{\hat{\pi}(\vect{X})(x_0) - \pi(x_0)}^2} \geq \frac{L^2}{64} \p{\frac{1}{4L^2}}^{2/3} n^{-2/3}
    \;.
\end{equation*}
This implies the first lower bound. 

\paragraph{$\epsilon$-DP overhead.}
By \Cref{eq:ghbbvcdjyuhsgdbefr} and by Le Cam's lemma for differential privacy on product distributions (\Cref{fact:lecamslemmaprivate}),
\begin{equation*}
    \begin{aligned}
    \inf_{\hat{\pi} \text{ $\epsilon$-DP}} \sup_{\pi \in \Theta^{\text{Lip}}_L}
    \E_{\vect{X} \sim \Prob_{\pi}^{\otimes n}, \hat{\pi}}\p{\p{\hat{\pi}(\vect{X})(x_0) - \pi(x_0)}^2} 
    &\geq 
    \frac{L^2}{8}h^2 (1-h)^2
    %(h^2-h)^2
    e^{-n\epsilon \tv{\Prob_{f_{L, x_0, h}}}{\Prob_{g}} } \\
    &\stackrel{\eqref{eq:jnbcdshgvdy}}{\geq} 
    \frac{L^2}{8} h^2 (1-h)^2 %(h^2-h)^2
    e^{-L n \epsilon h^2}\;.
    \end{aligned}
\end{equation*}
When possible (i.e. when $n \epsilon$ is big enough), setting $h = 1/\sqrt{n\epsilon}$ leads to, when $n \epsilon$ is large enough to ensure $1-h \geq 1/2$,  %$h - h^2 \geq h/2$,
\begin{equation*}
    \inf_{\hat{\pi} \text{ $\epsilon$-DP}} \sup_{\pi \in \Theta^{\text{Lip}}_L}
    \E_{\vect{X} \sim \Prob_{\pi}^{\otimes n}, \hat{\pi}}\p{\p{\hat{\pi}(\vect{X})(x_0) - \pi(x_0)}^2} \geq \frac{L^2 e^{-L}}{32} (n \epsilon)^{-1}
    \;.
\end{equation*}


\paragraph{$\rho$-zCDP overhead.}
By Le Cam's lemma for zero-concentrated differential privacy on product distributions (\Cref{fact:lecamslemmaconcentrated}) in \eqref{eq:ghbbvcdjyuhsgdbefr},
\begin{equation*}
    \begin{aligned}
    \inf_{\hat{\pi} \text{ $\epsilon$-DP}} \sup_{\pi \in \Theta^{\text{Lip}}_L}
    \E_{\vect{X} \sim \Prob_{\pi}^{\otimes n}, \hat{\pi}}\p{\p{\hat{\pi}(\vect{X})(x_0) - \pi(x_0)}^2} 
    &\geq 
    \frac{L^2}{8} h^2 (1-h)^2
    %(h^2-h)^2
    \p{1 - n \sqrt{\rho/2}\tv{\Prob_{f_{L, x_0, h}}}{\Prob_{g}}}  \\
    &\stackrel{\eqref{eq:jnbcdshgvdy}}{\geq} 
    \frac{L^2}{8}h^2 (1-h)^2 %(h^2-h)^2
    \p{1 - n \sqrt{\rho/2}Lh^2} \;.
    \end{aligned}
\end{equation*}
When possible (i.e. when $n \sqrt{\rho}$ is large enough), setting $h = \p{\frac{1}{\sqrt{2} L n \sqrt{\rho}}}^{1/2}$ leads to, when $n \sqrt{\rho}$ is large enough (so that $1-h \geq 1/2$ %$h - h^2 \geq h/2$
),
\begin{equation*}
    \inf_{\hat{\pi} \text{ $\rho$-zCDP}} \sup_{\pi \in \Theta^{\text{Lip}}_L}
    \E_{\vect{X} \sim \Prob_{\pi}^{\otimes n}, \hat{\pi}}\p{\p{\hat{\pi}(\vect{X})(x_0) - \pi(x_0)}^2} \geq \frac{L}{64} (n \sqrt{\rho})^{-1}
    \;.
\end{equation*}


\section{Proof of \Cref{th:lidp}}
\label{proofofth:lidp}

    Let $m \in \N \setminus \{0\}$ that will be fixed later in the proof. 
    As explained in the sketch of the proof, we build a packing consisting of functions that are parametrized by a vector $\omega \in \{ 0, 1\}^m$. After controlling quantities such as their pairwise TV distances, and their KL divergences to the uniform distribution, we leverage Fano-type inequalities in order to obtain lower-bounds.
    \paragraph{Packing construction.} 
    For any $\omega \in \{ 0, 1 \}^m$ different from $0$ and any $h>0$, we define the function $g_{L, \omega, h}$ as 
\begin{equation}
\label{eq:deftrianglemultiplefunction}
    g_{L, \omega, h} 
    \eqdef\frac{1}{\| \omega \|_1} \sum_{i=1}^m \omega_i f_{\| \omega \|_1 L, \frac{i}{m+1}, h} \;,
\end{equation}
where the functions $f$ are defined in \eqref{eq:deftrianglefunction}.
Note that $g_{L, \omega, h}$ is $L$-Lipschitz and that as soon as $h \leq h_m \eqdef \frac{1}{2(m+1)}$ it is also a valid density so that  $g_{L, \omega, h} \in \Theta^{\text{Lip}}_L$. Notice that the function $g_{L, \omega, h}$ is constant to $1-\| \omega \|_1 Lh^2$ everywhere except on each interval $\left[\frac{i}{m+1}-h, \frac{i}{m+1}+h\right]$ with $i$ such that $\omega_i \neq 0$, on which it deviates by a triangle of slopes $+L$ and $-L$. 

By denoting by $K$ the triangle kernel such that $K(t) = L t \Ind_{[-h, 0]}(t) - L t \Ind_{(0, h]}(t)$, it might be easier to visualize $g_{L, \omega, h}$ as
\begin{equation}
\label{nddgserrzw}
    \forall t \in [0, 1], \quad
    g_{L, \omega, h}(t) = 1 - \| \omega \|_1 \int K + \sum_{i=1}^m \omega_i K\p{t - \frac{1}{m+1}} \;,
\end{equation}
where $\int K = Lh^2$.

For $\omega, \omega' \in \{0, 1\}^m $ and for $h$ small enough (i.e. $h \leq h_m$), we can bound the total variation between $\Prob_{g_{L, \omega, h}}$ and $\Prob_{g_{L, \omega', h}}$ as 
% \begin{equation}
% \label{eq:tvsaw}
% \begin{aligned}
%     \tv{\Prob_{g_{L, \omega, h}}}{\Prob_{g_{L, \omega', h}}} 
%     &=
%     1 - \int_{[0, 1]} \min \p{g_{L, \omega, h}, g_{L, \omega', h}} \\
%     &\stackrel{\text{Constant part}}{\leq }
%     1 - \min \p{1-\| \omega \|_1 Lh^2, 1-\| \omega' \|_1 Lh^2} \\
%     &= \max \p{\|\omega\|_1 , \| \omega' \|_1 } Lh^2 \leq mLh^2\;.
% \end{aligned}
% \end{equation}
%\begin{equation}
\begin{align}
    \tv{\Prob_{g_{L, \omega, h}}}{\Prob_{g_{L, \omega', h}}} 
    &=
    \frac{1}{2} \int \left| g_{L, \omega, h} - g_{L, \omega', h} \right| \notag\\
    &\stackrel{\eqref{nddgserrzw}}{=}
    \frac{1}{2} \int \left| \| (\omega' \|_1 - \| \omega \|_1) \int K + \sum_{i=1}^m (\omega_i' - \omega_i) K\p{\cdot - \frac{1}{m+1}}\right| \notag\\
    &\leq
    \frac{1}{2} \int \left| \| \omega' \|_1 - \| \omega \|_1\right| \int K + \sum_{i=1}^m \left|\omega_i' - \omega_i\right| K\p{\cdot - \frac{1}{m+1}} \notag \\
    &=
    \frac{1}{2} \bigg( \left| \| \omega' \|_1 - \| \omega \|_1\right|  + \ham{\omega}{\omega'}\bigg) \int K \label{eq:tvsawham}\\
    &\leq mLh^2\;.\label{eq:tvsaw}\
\end{align}
%\end{equation}

The KL divergence between $\Prob_{g_{L, \omega, h}}$ and $\Prob_g$,
with $g$ the density constant equal to $1$ on $[0, 1]$, satisfies
\begin{equation}
\label{eq:klsaw}
    \begin{aligned}
    \kl{\Prob_{g_{L, \omega, h}}}{\Prob_g}
    &= 
    \int_{[0, 1]} \ln\p{g_{L, \omega, h}} g_{L, \omega, h} \\
    &=
    \ln \p{1 - \|\omega\|_1Lh^2} \p{1 - \|\omega\|_1Lh^2} \p{1 - \|\omega\|_1 2h } \\
    &\quad\quad\quad\quad+
    2 \|\omega\|_1 \int_{0}^h \ln \p{1 - \|\omega\|_1 Lh^2 + Lt} \p{1 - \|\omega\|_1 Lh^2 + Lt} dt \\
    &\stackrel{\ln (1 + \cdot) \leq \cdot}{\leq }
   \p{- \|\omega\|_1Lh^2} \p{1 -\|\omega\|_1 Lh^2} \p{1 - \|\omega\|_1 2h } \\
   &\quad\quad\quad\quad+
    2 \|\omega\|_1 \int_{0}^h \p{ - \|\omega\|_1 Lh^2 + Lt} \p{1 - \|\omega\|_1 Lh^2 + Lt} dt \\
    &\stackrel{\text{Calculus}}{=} \frac{L^2}{3} \|\omega\|_1 h^3 (2 - 3 \|\omega\|_1 h) \;.
    \end{aligned} 
\end{equation}
%\mRG{Expliquer la derniere ligne \CL{Fait} Je voulais dire macher un peu le travail pour voir comment mener les calculs en question. Deux lignes de plus ne feraient pas de mal !}

Finally, we lower bound the squared $L^2$ distance between $g_{L, \omega, h}$ and $g_{L, \omega', h}$:

\begin{equation}
\label{eq:kjuihqabzsdkjhfbqsd}
\begin{aligned}
    \int_{[0, 1]} &\p{g_{L, \omega, h} - g_{L, \omega', h}}^2 \\
    &= 
    \sum_{i=1}^m \Ind_{\omega_i \neq \omega_i'} \int_{\frac{i}{m+1} - h}^{\frac{i}{m+1} + h}
    \p{\p{\| \omega' \|_1 - \| \omega \|_1} \int K +(\omega_i-\omega'_i)  K\p{t-\frac{i}{m+1} }}^2 dt\\
    &\geq 
    \sum_{i=1}^m \Ind_{\omega_i \neq \omega_i'} \int_{\frac{i}{m+1} -h}^{\frac{i}{m+1} + h}
    \p{K\p{t-\frac{i}{m+1}} -  \left|\| \omega \|_1 - \| \omega' \|_1\right| \int K}^2 dt \\
    &\geq
    \sum_{i=1}^m \Ind_{\omega_i \neq \omega_i'} \int_{\frac{i}{m+1} - h}^{\frac{i}{m+1} + h}
    \left\{\p{K\p{t-\frac{i}{m+1}}}^2 \right.\\
    &\quad\quad\quad\quad\quad\quad\quad\quad\quad\quad \left.- 2 K\p{t-\frac{i}{m+1}} \left|\| \omega \|_1 - \| \omega' \|_1\right| \int K \right\} dt\\ 
    &\geq
    \ham{\omega}{\omega '} \p{  \int K^2 - 2m \p{\int K}^2} \\
    &\geq
    2 \ham{\omega}{\omega '} L^2 \p{  \frac{h^3}{3} - m h^4} \\
    &= \frac{2 \ham{\omega}{\omega'} L^2 h^3\p{1- 3 mh}}{3}\;.
\end{aligned}
\end{equation}
%The first inequality simply comes \mRG{JE trouve que ce serait plus clair en découpant en quelques lignes de calcul avec un explication simple pour chaque étape}from the fact when $\omega_i \neq \omega'_i$, a triangle is added around $\frac{i}{m+1}$ on one of the functions only. In the worst case (in terms of lower-bounds), it is added to the one that has the smallest constant part. On the support of the kernel, the integrated square difference between the two functions is then bigger than the same quantity integrated on the subset of the support where the function with the smallest continuous component is bigger than the other one.

By the Varshamov-Gilbert theorem \citep[Lemma 2.7]{tsybakov2003introduction}, as long as $m \geq 8$, there exist $M \in \N$ and $\omega^{(0)}, \dots, \omega^{(M)} \in \{ 0, 1 \}^m$ such that $M \geq 2^{m/8}$, $\omega^{(0)} = \{0\}^m$ and $i \neq j \implies \ham{\omega^{(i)}}{\omega^{(j)}} \geq m/8$. According to \eqref{eq:kjuihqabzsdkjhfbqsd}, the family $\p{g_{L, \omega^{(i)}, h}}_{i=1, \dots, M}$ is then %a
an $\Omega \coloneqq \frac{1}{2} \sqrt{\frac{mL^2\p{h^3 - 3 mh^4}}{12}}$ packing of $\Theta^{\text{Lip}}_L$ for the $L^2$ distance.

\paragraph{Recovering the usual lower-bound.}

By \Cref{jkhgqskdjhfbgkjqhwgsdkfjhgkqwsjhdgfjkhwsgdb} with $\Phi(\cdot) \coloneqq (\cdot)^2$ and $\|\cdot\|$ the $L^2$ norm,
\begin{equation}
\label{eq:dqslijuhfvsjkdcv}
    \begin{aligned}
    \inf_{\hat{\pi}  \text{ s.t. } \mathcal{C}} &\sup_{\pi \in \Theta^{\text{Lip}}_L}
    \E_{\vect{X} \sim \Prob_{\pi}^{\otimes n}, \hat{\pi}}\p{ \int_{[0, 1]}\p{\hat{\pi}(\vect{X})- \pi}^2} \\
    &\geq 
    \frac{mL^2 h^3\p{1 - 3 mh}}{48}
    \inf_{\hat{\pi}  \text{ s.t. } \mathcal{C}} \inf_{\Psi : \Theta^{\text{Lip}}_L \rightarrow \{0, 1\}}
    \max_{i = 1, \dots, M}  \Prob_{\vect{X} \sim \Prob_{g_{L, \omega^{(i)}, h}}^{\otimes n}, \hat{\pi}}                \p{\Psi(\hat{\pi}(\vect{X})) \neq i}  \\
    &\stackrel{\text{\Cref{fact:fanoslemma}}}{\geq}
    \frac{mL^2h^3\p{1 - 3 mh}}{48}
    \p{1 - \frac{1 + \frac{1}{M} \sum_{1\leq i\leq M} \kl{\Prob_{g_{L, \omega^{(i)}, h}}^{\otimes n}}{\Prob_{g}^{\otimes n}}}{\ln (M)}} \\
    &\stackrel{\text{Tensorization}}{=}
    \frac{mL^2 h^3\p{1 - 3 mh}}{48}
    \p{1 - \frac{1 + \frac{n}{M} \sum_{1\leq i\leq M} \kl{\Prob_{g_{L, \omega^{(i)}, h}}}{\Prob_{g}}}{\ln (M)}} \\
    &\stackrel{\eqref{eq:klsaw} \& \|\omega\|_1\leq m, M \geq 2^{m/8}}{\geq}
    \frac{mL^2 h^3\p{1 - 3 mh}}{48}
    \p{1 - \frac{1 +  \frac{L^2}{3}n m h^3 (2 - 3m h)}{\ln(2) m/8}} \;.
    \end{aligned}
\end{equation}
So, by choosing $m = \ceil{n^{
1/3}}$ and $h = \frac{c}{m}$ where c is a positive constant small enough we get, for $n$ big enough, 
\begin{equation*}
    \inf_{\hat{\pi} \text{ $\epsilon$-DP}} \sup_{\pi \in \Theta^{\text{Lip}}_L}
    \E_{\vect{X} \sim \Prob_{\pi}^{\otimes n}, \hat{\pi}}\p{ \int_{[0, 1]}\p{\hat{\pi}(\vect{X})- \pi}^2} \geq C^{-1} (n)^{-2/3} \;,
\end{equation*}
where $C$ is a positive constant depending only on $L$.

\paragraph{$\epsilon$-DP overhead.}

By the same reduction and Fano's lemma for differential privacy on product distributions (\Cref{fact:fanoslemmaprivate}), we get for any $h \leq h_m$,
\begin{equation*}
    \begin{aligned}
    \inf_{\hat{\pi} \text{ $\epsilon$-DP}} \sup_{\pi \in \Theta^{\text{Lip}}_L}
    &\E_{\vect{X} \sim \Prob_{\pi}^{\otimes n}, \hat{\pi}}\p{ \int_{[0, 1]}\p{\hat{\pi}(\vect{X})- \pi}^2} \\
    &\geq 
    \frac{mL^2 h^3\p{1 - 3 mh}}{48}
    \p{1 - \frac{1 + \frac{n \epsilon}{M^2} 2 \sum_{1\leq i, j \leq M} \tv{\Prob_{g_{L, \omega^{(i)}, h}}}{\Prob_{g_{L, \omega^{(j)}, h}}}}{\ln (M)}} \\
    &\stackrel{\eqref{eq:tvsaw} \& M\geq 2^{m/8}}{\geq }
    \frac{mL^2h^3  \p{1- 3 mh}}{48}
    \p{1 - \frac{1 + 2 n \epsilon m L h^2}{\ln(2) m/8}} \;.
    \end{aligned}
\end{equation*}
So, by choosing $m = \ceil{\sqrt{n \epsilon}}$ and $h = \frac{c}{m}$  where c is small enough a positive constant (depending only on $L$), we get, %for $n$ and $n \epsilon$ big enough, 
as soon as $\min(n,n\epsilon)$ is big enough,
\begin{equation*}
    \inf_{\hat{\pi} \text{ $\epsilon$-DP}} \sup_{\pi \in \Theta^{\text{Lip}}_L}
    \E_{\vect{X} \sim \Prob_{\pi}^{\otimes n}, \hat{\pi}}\p{ \int_{[0, 1]}\p{\hat{\pi}(\vect{X})- \pi}^2} \geq C'^{-1} (n\epsilon)^{-1} \;,
\end{equation*}
where $C'$ is a positive constant depending only on $L$.

\paragraph{$\rho$-zCDP overhead.}For $\rho$-zCDP, we present the proof using both Fano's lemma and Assouad's method. We will see that Assouad gives better results

\paragraph{Fano version.}
By again the same reduction and Fano's lemma for zero-concentrated differential privacy (\Cref{fact:fanoslemmaconcentrated}), denoting $t_{i, j} \eqdef \tv{\Prob_{g_{L, \omega^{(i)}, h}}}{\Prob_{g_{L, \omega^{(j)}, h}}}$, we get for any $h \leq h_m$,
\begin{equation*}
    \begin{aligned}
    \inf_{\hat{\pi} \text{ $\rho$-zCDP}} \sup_{\pi \in \Theta^{\text{Lip}}_L}
    &\E_{\vect{X} \sim \Prob_{\pi}^{\otimes n}, \hat{\pi}}\p{ \int_{[0, 1]}\p{\hat{\pi}(\vect{X})- \pi}^2} \\
    &\geq 
    \frac{mL^2h^3\p{1 - 3 mh}}{48}
    \p{1 - \frac{1 + \frac{n^2 \rho}{M^2} 4 \sum_{1\leq i, j \leq M} \p{\frac{1}{n} t_{i, j} + t_{i, j}^2}}{\ln (M)}}  \\
    &\stackrel{\eqref{eq:tvsaw}}{\geq }
    \frac{mL^2h^3\p{1 - 3 mh}}{48}
    \p{1 - \frac{1 + n^2 \rho 4 \p{\frac{mLh^2}{n} + m^2 L^2 h^4}}{\ln (2) m / 8}} 
     \;.
    \end{aligned}
\end{equation*}
So, by choosing $m = \ceil{\p{n \sqrt{\rho}}^{
\frac{2}{3}}}$ and $h = \frac{c}{m} $ for $c$ small enough (depending only on $L$), if $\frac{n}{\rho}$ is big enough, we get that 
\begin{equation*}
    \inf_{\hat{\pi} \text{ s.t. } \text{$\rho$-zCDP}} \sup_{\pi \in \Theta^{\text{Lip}}_L}
    \E_{\vect{X} \sim \Prob_{\pi}^{\otimes n}, \hat{\pi}}\p{ \int_{[0, 1]}\p{\hat{\pi}(\vect{X})- \pi}^2} \geq C''^{-1} (n \sqrt{\rho})^{-4/3}
\end{equation*}
where $C''$ is a positive constant depending only on $L$.

\paragraph{Assouad version.}

From \Cref{eq:kjuihqabzsdkjhfbqsd}, we can see that when $h \eqdef \frac{c}{m}$ for a positive $c$ that is small enough, the condition expressed in \Cref{eq:Assouadcondition} is satisfied for $\tau = \Omega(h^{3})$. To apply~\eqref{eq:AssouadMixtures}, the only missing ingredient is to bound the testing difficulties between the mixtures on the hypercube.

In the sequel, $\Prob_{\omega}$ is used as a short for $\Prob_{g_{L, \omega, h}}$. We need to bound the total variation between the mixtures on the hypercube (see~\eqref{eq:AssouadMixtures}) as%\RG{Deuxieme ligne, pourquoi $1/2^m$ et pas $1/2^{m-1}$?}
\begin{equation*}
\begin{aligned}
    \tv{\Prob_{\omega^{i,0}}}{\Prob_{\omega^{i,1}}}
     &= \tv{\frac{1}{2^{m-1}} \sum_{\omega \in \{0, 1 \}^m | \omega_i = 0} \Prob_{g_{L, \omega, h}} }{\frac{1}{2^{m-1}} \sum_{\omega \in \{0, 1 \}^m | \omega_i = 1} \Prob_{g_{L, \omega, h}} } \\
    &= \frac{1}{2} \frac{1}{2^{m-1}} \int \left| \sum_{\omega \in \{0, 1 \}^m | \omega_i = 0} g_{L, \omega, h} - \sum_{\omega \in \{0, 1 \}^m | \omega_i = 1} g_{L, \omega, h} \right| \\
    &= \frac{1}{2^m} \int \Bigg| \sum_{\omega_1, \dots, \omega_{i-1}, \omega_{i+1} \dots, \omega_{m} \in \{0, 1 \}} \big( g_{L,  (\omega_1, \dots, \omega_{i-1}, 0, \omega_{i+1} \dots, \omega_{m}), h} - \\  
    &\quad\quad\quad\quad\quad\quad\quad\quad\quad\quad\quad\quad\quad\quad\quad\quad\quad\quad\quad\quad g_{L, (\omega_1, \dots, \omega_{i-1}, 1, \omega_{i+1} \dots, \omega_{m}), h} \big) \Bigg| \\
    &\leq \frac{1}{2^m} \sum_{\omega_1, \dots, \omega_{i-1}, \omega_{i+1} \dots, \omega_{m} \in \{0, 1 \}} \int \Bigg|   g_{L, (\omega_1, \dots, \omega_{i-1}, 0, \omega_{i+1} \dots, \omega_{m}), h} - \\  
    &\quad\quad\quad\quad\quad\quad\quad\quad\quad\quad\quad\quad\quad\quad\quad\quad\quad\quad\quad\quad g_{L,  (\omega_1, \dots, \omega_{i-1}, 1, \omega_{i+1} \dots, \omega_{m}), h}  \Bigg| \\
    &\stackrel{\eqref{eq:tvsawham}}{\leq} \frac{1}{2^m} \sum_{\omega_1, \dots, \omega_{i-1}, \omega_{i+1} \dots, \omega_{m} \in \{0, 1 \}} 2 L h^2  \\
    &= O \p{h^{2}}\;.
\end{aligned}  
\end{equation*}
%Now, we can notice that in this sum, $(\omega_1, \dots, \omega_{i-1}, 0, \omega_{i+1} \dots, \omega_{m})$ and $(\omega_1, \dots, \omega_{i-1}, 1, \omega_{i+1} \dots, \omega_{m})$ are always at Hamming distance $1$. 
%The difference $g_{L, (\omega_1, \dots, \omega_{i-1}, 0, \omega_{i+1} \dots, \omega_{m}), h} - g_{L, (\omega_1, \dots, \omega_{i-1}, 1, \omega_{i+1} \dots, \omega_{m}), h}$ can thus be simplified a lot \RG{cad? Serait-ce bcp plus long d'écrire explicitement pour éviter au lecteur / reviewer de multiples aller-retours ?}. By %plugin
%\RG{plugging} back the definition \Cref{eq:deftrianglemultiplefunction} , and by the triangular inequality, by denoting by $K$ the triangle kernel such that $K(t) = L t \Ind_{[-h, 0]}(t) - L t \Ind_{(0, h]}(t)$, 
%\RG{Il me semble que tout serait bien plus simple si en fait tu commençais par définir le noyau triangulaire $K_{L,h}(t)$, et que tu l'utilisais pour définir $f_{L,x,h}$ et/ou $g_{L,\omega,h}$.}
%we get
%\begin{equation*}
%\begin{aligned}
%\tv{\Prob_{\omega^{i,0}}}{\Prob_{\omega^{i,1}}}
%    &\tv{\frac{1}{2^{m-1}} \sum_{\omega \in \{0, 1 \}^m | \omega_i = 0} \Prob_{g_{L, \omega, h}} }{\frac{1}{2^{m-1}} \sum_{\omega \in \{0, 1 \}^m | \omega_i = 1} \Prob_{g_{L, \omega, h}} } \\
%    &\leq \frac{1}{2^m} \sum_{\omega_1, \dots, \omega_{i-1}, \omega_{i+1} \dots, \omega_{m} \in \{0, 1 \}} \int \Bigg( L h^2 + K  \Bigg) \\
%    &\leq \frac{1}{2^m} \sum_{\omega_1, \dots, \omega_{i-1}, \omega_{i+1} \dots, \omega_{m} \in \{0, 1 \}} 2 L h^2  \\
%    &= O \p{h^{2}}\;.
%\end{aligned}  
%\end{equation*}
%\RG{Tel quel sans borne de l'intégrale ça parait faux}
%\RG{Qu'est-ce que $K$ dans la premiere inégalité ? Il manque bcp de détails. 
%\RG{Ne faudrait-il pas prouver explicitement un équivalent de \eqref{eq:tvsaw}? }
%\CL{J'ai ajouté du détail sur ce qu'est $K$}}
Here and in the sequel, the asymptotic comparators only hide constants and terms that depend on $L$.
All in all, by using Le Cam's lemma for product distribution and $\rho$-zCDP \Cref{fact:lecamslemmaconcentrated}, and by \Cref{eq:Assouadprivate}, since $\tau = \Omega(h^3)$ we obtain
\begin{equation}
    \inf_{\hat{\pi} \text{ $\rho$-zCDP}} \sup_{\pi \in  \Theta^{\text{PSob}}_{L, \beta}}
    \E_{\vect{X} \sim \Prob_{\pi}^{\otimes n}, \hat{\pi}}\p{ \int_{[0, 1]}\p{\hat{\pi}(\vect{X})- \pi}^2} = \Omega \p{m h^{3}} \p{1 - n \sqrt{\rho} O \p{h^{2}}}.
\end{equation}
Setting $h \approx \p{n \sqrt{\rho}}^{\frac{-1}{2}}$ concludes the proof.
    

\section{proof of \Cref{lemma:projectionutility}}
\label{proofoflemma:projectionutility}


Let $\pi \in \Theta^{\text{PSob}}_{L, \beta}$.
    We have, 
    \begin{equation*}
    \begin{aligned}
         \E \p{ \int_{[0, 1]} \p{\hat{\pi}^{\text{proj}}(\vect{X}) - \pi}^2}
         & \stackrel{\text{Parseval}}{=}
          \E \p{ \sum_{i=1}^N \p{\hat{\theta}_i -\theta_i + \frac{1}{n} Z_i}^2 + \sum_{i=N+1}^{+\infty} \theta_i^2 } \\
          & \stackrel{\text{}}{=}
           \sum_{i=1}^N \E \p{\p{\hat{\theta}_i -\theta_i + \frac{1}{n} Z_i}^2} + \sum_{i=N+1}^{+\infty} \theta_i^2 \;.
    \end{aligned}
    \end{equation*}

Furthermore, for any $i$, since $Z$ is centered
\begin{equation*}
    \begin{aligned}
        \E \p{\hat{\theta}_i} 
        &= 
        \E \p{\frac{1}{n} \sum_{j=1}^n \phi_i(X_j)}
        =
        \frac{1}{n} \sum_{j=1}^n \E \p{\phi_i(X_j)} 
        \stackrel{X_j\ \mathrm{i.i.d.}}{=} \E_{\vect{X} \sim \Prob_{\pi}^{\otimes n}} \phi_i(X_1) = \int \pi \phi_i = \theta_i
        % \\
        % &\stackrel{}{=}
        % \frac{1}{n} \sum_{j=1}^n \int_{[0, 1]} \dots \int_{[0, 1]} \phi_i(x_j) \pi(x_1) \dots \pi(x_n) dx_1 \dots dx_n \\
        % &\stackrel{}{=}
        % \frac{1}{n} \sum_{j=1}^n \int_{[0, 1]} \phi_i(x_j) \pi(x_j) dx_j 
        % = \frac{1}{n} \sum_{j=1}^n \theta_i = \theta_i \;.
    \end{aligned}
\end{equation*}
Hence, for any $i$, since $Z_i$ is independent from the dataset
\begin{equation*}
    \begin{aligned}
        \E \p{\p{\hat{\theta}_i -\theta_i + \frac{1}{n} Z_i}^2}
        &= 
        \V \p{\hat{\theta}_i} + \frac{1}{n^2} \V \p{Z_i} 
        \stackrel{\text{Independence of $X_j$}}{=}
        \frac{1}{n^2} \sum_{j=1}^n \V \p{\phi_i \p{X_j}} + \frac{1}{n^2} \V \p{Z_i} \\
        &\stackrel{|\phi_i| \leq \sqrt{2}}{\leq}
        \frac{%8
        2}{n}+ \frac{1}{n^2} \V \p{Z} \;.
    \end{aligned}
\end{equation*}
Finally, with $a_j \coloneqq j-1$, \Cref{lemma:sobolevellipsoid} allows bounding $\sum_{i=m+1}^{+\infty} \theta_i^2$ as
\begin{equation*}
    \sum_{i=N+1}^{+\infty} \theta_i^2 \leq \frac{1}{N^{2 \beta}} \sum_{i=N+1}^{+\infty} a_i^{2\beta} \theta_i^2 \leq \frac{1}{N^{2 \beta}} \sum_{i=1}^{+\infty} a_i^{2\beta} \theta_i^2 \stackrel{\Cref{lemma:sobolevellipsoid}}{\leq} \frac{1}{N^{2 \beta}} \frac{L^2}{\pi^{2 \beta}} \;.
\end{equation*}
This yields the conclusion with $C_{L,\beta} \coloneqq \max(2, L^2/\pi^{2\beta})$.

\section{Proof of  \Cref{th:sidp}}
\label{proofofth:sidp}

Let us consider the following well-known function :
\begin{equation*}
    \forall x \in \R, \quad K_0(x) \eqdef e^{- \frac{1}{1 - x^2}}\Ind_{(-1, 1)}(x) \;.
\end{equation*}
We can notice that for any $\beta>0$ there exists $\nu > 0$ such that the kernel $K(x) \eqdef \nu K_0(2x)$ satisfies $K \in \mathcal{C}^{\infty}(\R, [0, +\infty))$, $\int \p{K^{(\beta)}}^2 \leq 1$  %\mRG{A-t-on défini la notation $K^{(\beta)}$ qq part ?} 
and $K(x)>0$ iff $x \in (-1/2, 1/2)$. Furthermore, for any $i \in \N$, $K^{(i)}(x) = 0$ for every $x \in (-\infty, -1/2] \cup [1/2, +\infty)$.

\paragraph{Packing construction.}
%\mRG{D'apres ce que je sais on écrit soit "Let $m$ {\bf be} an integer ..." soit "Consider $m$ an integer}
Let $m \in \N \setminus \{0\}$ that will be fixed later. For any $h > 0$, and $\omega \in \{0, 1\}^m$, we define the function $g_{L, \beta, \omega, h}$ as,
 \begin{equation}
 \label{eq:defgkernel}
     \forall x \in [0, 1],
     \quad 
     g_{L, \beta, \omega, h}(x) \coloneqq 1 - \| \omega \|_1 L h^{\beta+1} \int K + L h^{\beta} \sum_{i=1}^m \omega_i K\p{\frac{x - \frac{i}{m+1}}{h}} \;.
\end{equation}
Note that when $h < \frac{1}{m+1}$ we have $\int_0^1 g_{L,\beta,\omega,h} =1$; when $m h \int \p{K^{(\beta)}}^2 \leq 1$, we have $g_{L,\beta,\omega,h} \geq 0$; and when both hold we have $g_{L, \beta, \omega, h} \in \Theta^{\text{PSob}}_{L, \beta}$ (see \Cref{eq:defpsob}). %\RG{une ligne pour expliquer pourquoi $\int (g^{(\beta)})^2 \leq L^2$ pour éviter au lecteur de devoir prendre une feuille de papier entre chaque ligne ?}
Indeed, under these hypotheses, the periodicity conditions are immediate (the function is constant on neighborhoods of $0$ and $1$, with the same value). The energy of the $\beta$th derivative can be bounded as
\begin{equation*}
\begin{aligned}
    \int \p{g_{L, \beta, \omega, h}^{(\beta)}}^2
    &= \int \p{L h^{\beta} \sum_{i=1}^m \omega_i \p{x \mapsto K\p{ \frac{x - \frac{i}{m+1}}{h}}}^{(\beta)}}^2 \\
    &= \int \p{L\sum_{i=1}^m \omega_i K^{(\beta)}\p{ \frac{\cdot - \frac{i}{m+1}}{h}}}^2 \\
    &\stackrel{\text{disjoint support}}{=} L^2 \sum_{i=1}^m \omega_i\int  \p{K^{(\beta)}\p{ \frac{\cdot - \frac{i}{m+1}}{h}}}^2 \\
    &= L^2 m h \int \p{K^{(\beta)}}^2 \leq L^2 \;.
\end{aligned}
\end{equation*}

In the sequel of this proof, this hypothesis will always be satisfied asymptotically for all the values of $m$ and $h$ that will be considered. From now on, we may consider it valid.

Given $h>0$ and $\omega, \omega' \in \{0, 1\}^m$, when $g_{L, \beta, \omega, h}, g_{L, \beta, \omega', h} \in \Theta^{\text{PSob}}_{L, \beta}$, we can bound the total variation between $\Prob_{g_{L, \beta, \omega, h}}$ and  $\Prob_{g_{L, \beta, \omega', h}}$ as,
%\begin{equation}
%\label{eq:tvkernelsaw}
%\begin{aligned}
%    \tv{\Prob_{g_{L, \beta, \omega, h}}}{\Prob_{g_{L, \beta, \omega', h}}} 
%    &= 
%    1 - \int_{[0, 1]} \min \p{g_{L, \beta, \omega, h}, g_{L, \beta, \omega', h}}\\
%    &\stackrel{\text{Constant part}}{\leq }
%    1 -  \p{1 - \max( \| \omega \|_1,  \| \omega' \|_1)L h^{\beta+1} \int K} \\
%    &= \max( \| \omega \|_1,  \| \omega' \|_1) L h^{\beta+1} \int K \leq m L h^{\beta+1} \int K \;.
%\end{aligned}
%\end{equation}
%\notag

\begin{align}
    \tv{\Prob_{g_{L, \beta, \omega, h}}}{\Prob_{g_{L, \beta, \omega', h}}} 
    &=
    \frac{1}{2} \int \left| g_{L, \beta, \omega, h} - g_{L, \beta, \omega', h} \right| \notag \\
    &=
    \frac{1}{2} \int \left| \| (\omega' \|_1 - \| \omega \|_1) L h^{\beta + 1} \int K + \sum_{i=1}^m (\omega_i' - \omega_i) L h^{\beta } K\p{\frac{\cdot - \frac{1}{m+1}}{h}}\right| \notag\\
    &\leq
   \frac{1}{2} \int \| |\omega' \|_1 - \| \omega \|_1| L h^{\beta + 1} \int K + \sum_{i=1}^m |\omega_i' - \omega_i| L h^{\beta } K\p{\frac{\cdot - \frac{1}{m+1}}{h}} \notag\\
    &=
    \frac{1}{2} \bigg( \left| \| \omega' \|_1 - \| \omega \|_1\right|  + \ham{\omega}{\omega'}\bigg) L h^{\beta+1} \int K \label{eq:tvkernelsawham}\\
    &\leq mL h^{\beta+1} \label{eq:tvkernelsaw}\;.
\end{align}

The KL divergence between $\Prob_{g_{L, \beta, \omega, h}}$ and $\Prob_{g}$, the uniform distribution on $[0, 1]$, is bounded as
\begin{equation}
    \label{eq:klkernelsaw}
    \begin{aligned}
        &\kl{\Prob_{g_{L, \beta, \omega, h}}}{\Prob_{g}}
        =
        \int_{[0, 1]} \ln \p{g_{L, \beta, \omega, h}} g_{L, \beta, \omega, h} \\
        &=
        \int_{[0, 1] \setminus \cup_{i : \omega_i \neq 0} \left[ \frac{i}{m+1} - \frac{h}{2}, \frac{i}{m+1} + \frac{h}{2} \right]} 
        \ln \p{1 - \| \omega \|_1 L h^{\beta+1} \int K} \p{1 - \| \omega \|_1 L h^{\beta+1} \int K} dt \\ 
        &\quad\quad\quad\quad+ 
        \| \omega \|_1 \int_{-\frac{h}{2}}^{\frac{h}{2}} \ln \p{1 - \| \omega \|_1 L h^{\beta+1} \int K + L h^{\beta} K \p{\frac{t}{h}}}  \\
        &\quad\quad\quad\quad\quad\quad\quad\quad\quad\quad\quad\quad
        \p{1 - \| \omega \|_1 L h^{\beta+1} \int K + L h^{\beta} K \p{\frac{t}{h}}}dt \\
        &\stackrel{\ln (1 + \cdot) \leq \cdot}{\leq }
        \p{1 - \| \omega \|_1 h}
        \p{- \| \omega \|_1 L h^{\beta+1} \int K} \p{1 - \| \omega \|_1 L h^{\beta+1} \int K}  \\ 
        &\quad\quad\quad\quad+ 
        \| \omega \|_1 \int_{-\frac{h}{2}}^{\frac{h}{2}} \p{- \| \omega \|_1 L h^{\beta+1} \int K + L h^{\beta} K \p{\frac{t}{h}}}  \\
        &\quad\quad\quad\quad\quad\quad\quad\quad\quad\quad\quad\quad
        \p{1 - \| \omega \|_1 L h^{\beta+1} \int K + L h^{\beta} K \p{\frac{t}{h}}}dt\\
        &\stackrel{\text{Calculus}}{=}
        \| \omega \|_1 L^2 h^{2\beta + 1} \int K^2 - \| \omega \|_1^2 L^2 h^{2\beta + 2} \int K \\
        &\stackrel{}{\leq}
        \| \omega \|_1 L^2 h^{2\beta + 1} \int K^2
        \leq m L^2 h^{2\beta + 1} \int K^2 \;.
    \end{aligned}
\end{equation}
%\RG{Ca ne mangerait pas de pain de donner qq étapes du "Calculus" pour faciliter la vie au pauvre lecteur qui n'a pas de papier sous la main. Personnellement je lis entre 2 exposés et je n'ai pas vérifié, c'est dommage car le reste se suite bien}
Finally, the squared $L^2$ distance between $g_{L, \beta, \omega, h}$ and $g_{L, \beta, \omega', h}$ can be lower bounded as,
\begin{equation}
\label{eq:distkernelsaw}
\begin{aligned}
    \int_{[0, 1]} &\p{g_{L, \beta, \omega, h} - g_{L, \beta, \omega', h}}^2 \\
    &= 
    \sum_{i=1}^m \Ind_{\omega_i \neq \omega_i'} \int_{\frac{i}{m+1} - \frac{h}{2}}^{\frac{i}{m+1} + \frac{h}{2}}
    \p{L h^{\beta + 1} \p{\| \omega' \|_1 - \| \omega \|_1} \int K +(\omega_i-\omega'_i) L h^{\beta } K\p{\frac{t-\frac{i}{m+1} }{h}}}^2 dt\\
    &\geq 
    \sum_{i=1}^m \Ind_{\omega_i \neq \omega_i'} \int_{\frac{i}{m+1} - \frac{h}{2}}^{\frac{i}{m+1} + \frac{h}{2}}
    \p{L h^{\beta } K\p{\frac{t-\frac{i}{m+1}}{h}} - L h^{\beta + 1} \left|\| \omega \|_1 - \| \omega' \|_1\right| \int K}^2 dt \\
    &\geq
    \sum_{i=1}^m \Ind_{\omega_i \neq \omega_i'} \int_{\frac{i}{m+1} - \frac{h}{2}}^{\frac{i}{m+1} + \frac{h}{2}}
    \left\{\p{L h^{\beta } K\p{\frac{t-\frac{i}{m+1}}{h}}}^2 \right.\\
    &\quad\quad\quad\quad\quad\quad\quad\quad\quad\quad \left.- 2 L h^{\beta } K\p{\frac{t-\frac{i}{m+1}}{h}} L h^{\beta + 1} \left|\| \omega \|_1 - \| \omega' \|_1\right| \int K \right\} dt\\ 
    &\geq
    \ham{\omega}{\omega '} L^2h^{2 \beta + 1}\p{  \int K^2 - 2m h \p{\int K}^2} \;.
\end{aligned}
\end{equation}
%\RG{toutes les étapes sont super claires, s'en inspirer pour détailler \eqref{eq:kjuihqabzsdkjhfbqsd}?}
By the Varshamov-Gilbert theorem \citep[Lemma 2.7]{tsybakov2003introduction}, as long as $m \geq 8$, there exist $M \in \N$ and $\omega^{(0)}, \dots, \omega^{(M)} \in \{ 0, 1 \}^m$ such that $M \geq 2^{m/8}$, $\omega^{(0)} = \{0\}^m$ and $i \neq j \implies \ham{\omega^{(i)}}{\omega^{(j)}} \geq m/8$. According to \eqref{eq:distkernelsaw}, the family $\p{g_{L, \beta, \omega^{(i)}, h}}_{i=1, \dots, M}$ is then a $\Omega = \frac{1}{2} \sqrt{\frac{m}{8}  L^2 h^{2 \beta + 1} \p{ \int K^2 - 2m h \p{\int K}^2}}$ packing of $\Theta^{\text{PSob}}_{L, \beta}$ for the $L^2$ distance.

\paragraph{Recovering the usual lower-bound.}

By \Cref{jkhgqskdjhfbgkjqhwgsdkfjhgkqwsjhdgfjkhwsgdb} with $\Phi(\cdot) \coloneqq (\cdot)^2$ and $\|\cdot\|$ the $L^2$ norm,
\begin{equation}
\label{eq:dqslijuhfvsjkdcv}
    \begin{aligned}
    \inf_{\hat{\pi}  \text{ s.t. } \mathcal{C}} &\sup_{\pi \in \Theta^{\text{PSob}}_{L, \beta}}
    \E_{\vect{X} \sim \Prob_{\pi}^{\otimes n}, \hat{\pi}} \int_{[0, 1]}\p{\hat{\pi}(\vect{X})- \pi}^2 \\
    &\geq 
    \frac{L^2}{32}  m h^{2 \beta + 1} \p{  \int K^2 - 2m h \p{\int K}^2} \\
    &\quad\quad\quad\quad\inf_{\hat{\pi} \text{ s.t. } \mathcal{C}} \inf_{\Psi :\Theta^{\text{PSob}}_{L, \beta} \rightarrow \{0, 1\}}
    \max_{i = 1, \dots, M}   \Prob_{\vect{X} \sim \Prob_{g_{L, \beta, \omega^{(i)}, h}}^{\otimes n}, \hat{\pi}}                \p{\Psi(\hat{\pi}(\vect{X})) \neq i} \\
    &\stackrel{\text{\Cref{fact:fanoslemma}}}{\geq}
    \frac{L^2}{32} m h^{2 \beta + 1}  \p{ \int K^2 - 2m h \p{\int K}^2} \\
    &\quad\quad\quad\quad\quad\quad\quad\quad\quad\quad\quad\quad\quad\quad\quad\quad\quad\quad
    \p{1 - \frac{1 + \frac{1}{M} \sum_{1\leq i\leq M} \kl{\Prob_{g_{L, \beta, \omega^{(i)}, h}}^{\otimes n}}{\Prob_{g}^{\otimes n}}}{\ln (M)}} \\
    &\stackrel{\text{Tensorization}}{=}
   \frac{L^2}{32} mh^{2 \beta + 1} \p{  \int K^2 - 2m h \p{\int K}^2} \\
    &\quad\quad\quad\quad\quad\quad\quad\quad\quad\quad\quad\quad\quad\quad\quad\quad\quad\quad\quad\quad
    \p{1 - \frac{1 + \frac{n}{M} \sum_{1\leq i\leq M} \kl{\Prob_{g_{L, \beta, \omega^{(i)}, h}}}{\Prob_{g}}}{\ln (M)}} \\ 
    &\stackrel{\eqref{eq:klkernelsaw}\& M \geq 2^{m/8}}{\geq}
    \frac{L^2}{32} m h^{2 \beta + 1}  \p{ \int K^2 - 2m h \p{\int K}^2}
    \p{1 - \frac{1 + n m L^2 h^{2\beta + 1} \int K^2}{\ln (2) m/8}}\;.
    \end{aligned}
\end{equation}
%\RG{D'où sort le facteur $1/8$ au numérateur sous $\int K^2$? \CL{Typo en lien avec la ligne d'en dessous}}
Finally, setting $m = \ceil{ n^{\frac{1}{2 \beta + 1}}}$ and $h = \frac{c}{m}$ for $c$ small enough gives that, for $n$ big enough, 
\begin{equation*}
    \inf_{\hat{\pi} \text{ $\epsilon$-DP}} \sup_{\pi \in \Theta^{\text{PSob}}_{L, \beta}}
    \E_{\vect{X} \sim \Prob_{\pi}^{\otimes n}, \hat{\pi}} \int_{[0, 1]}\p{\hat{\pi}(\vect{X})- \pi}^2 \geq C^{-1} n^{-\frac{2 \beta}{2 \beta + 1}} \;,
\end{equation*}
where $C$ is a positive constant depending only on $L$ and $\beta$.

\paragraph{$\epsilon$-DP overhead.}

By the same reduction and Fano's lemma for differential privacy on product distributions (\Cref{fact:fanoslemmaprivate}), we get
\begin{equation*}
    \begin{aligned}
    \inf_{\hat{\pi} \text{ $\epsilon$-DP}} \sup_{\pi \in \Theta^{\text{PSob}}_{L, \beta}}
    &\E_{\vect{X} \sim \Prob_{\pi}^{\otimes n}, \hat{\pi}} \int_{[0, 1]}\p{\hat{\pi}(\vect{X})- \pi}^2 \\
    &\geq 
     \frac{L^2}{32}  m h^{2 \beta + 1}  \p{ \int K^2 - 2m h \p{\int K}^2} \\
    &\quad\quad\quad\quad \p{1 - \frac{1 + \frac{n \epsilon}{M^2} 2 \sum_{1\leq i, j \leq M} \tv{\Prob_{g_{L, \beta, \omega^{(i)}, h}}}{\Prob_{g_{L, \beta, \omega^{(j)}, h}}}}{\ln (M)}} \\
    &\stackrel{\eqref{eq:tvkernelsaw}}{\geq }
    \frac{L^2}{32}  m h^{2 \beta + 1} \p{  \int K^2 - 2m h \p{\int K}^2}
    \p{1 - \frac{1 + 2 n \epsilon m L h^{\beta+1} \int K}{\ln(2) m/8}} \;.
    \end{aligned}
\end{equation*}
Setting $m = \ceil{\p{n \epsilon}^{\frac{1}{\beta + 1}}}$ and $h = \frac{c}{m}$ for $c$ small enough leads to, for $n\epsilon$ big enough, %\RG{Ne suffit-il pas que $n\epsilon$ soit suffisamment grand ? Pourquoi faut-il aussi que $n$ soit grand?}
\begin{equation*}
    \begin{aligned}
    \inf_{\hat{\pi} \text{ $\epsilon$-DP}} \sup_{\pi \in \Theta^{\text{PSob}}_{L, \beta}}
    \E_{\vect{X} \sim \Prob_{\pi}^{\otimes n}, \hat{\pi}} \int_{[0, 1]}\p{\hat{\pi}(\vect{X})- \pi}^2 
    \geq 
     C'^{-1} \p{n \epsilon}^{-\frac{2 \beta}{\beta + 1}} \;,
    \end{aligned}
\end{equation*}
where $C'$ is a constant depending only on $L$ and $\beta$.

\paragraph{$\rho$-zCDP overhead.}For $\rho$-zCDP, we present the proof using both Fano's lemma and Assouad's method. We will see that Assouad gives better results.

\paragraph{Fano version.}
By again the same reduction and Fano's lemma for zero-concentrated differential privacy (\Cref{fact:fanoslemmaconcentrated}), denoting $t_{i, j} \eqdef \tv{\Prob_{g_{L, \beta, \omega^{(i)}, h}}}{\Prob_{g_{L, \beta,\omega^{(j)}, h}}}$, we get
%\RG{C'est au moins la troisième fois qu'on écrit quasiment les memes premieres lignes. Il y a matière à factoriser pas mal, non ?}
\begin{equation*}
    \begin{aligned}
    \inf_{\hat{\pi} \text{ $\rho$-zCDP}} \sup_{\pi \in  \Theta^{\text{PSob}}_{L, \beta}}
    &\E_{\vect{X} \sim \Prob_{\pi}^{\otimes n}, \hat{\pi}} \int_{[0, 1]}\p{\hat{\pi}(\vect{X})- \pi}^2 \\
    &\geq 
     \frac{L^2}{32}  m h^{2 \beta + 1} \p{  \int K^2 - 2m h \p{\int K}^2} \\
     &\quad\quad\quad\quad\quad\quad\quad\quad\quad\quad\quad\quad\quad\quad\quad\quad
    \p{1 - \frac{1 + \frac{n^2 \rho}{M^2} 4 \sum_{1\leq i, j \leq M} \frac{1}{n} t_{i, j} + t_{i, j}^2}{\ln (M)}}  \\
    &\stackrel{\eqref{eq:tvkernelsaw}}{\geq}
    \frac{L^2}{32}  m h^{2 \beta + 1}  \p{ \int K^2 - 2m h \p{\int K}^2} \\
    &\quad\quad\quad\quad\quad\quad\quad\quad\quad\quad\quad\quad
    \p{1 - \frac{1 + 4 n^2 \rho \p{\frac{m L h^{\beta+1} \int K}{n} +\p{m L h^{\beta+1} \int K}^2}}{\ln (2) m / 8}} 
     \;.
    \end{aligned}
\end{equation*}

So, by choosing $m =  \ceil{\p{n \sqrt{\rho}}^{\frac{2}{2 \beta + 1}}}$ and $h = \frac{c}{m}$ for $c$ small enough, if 
%\RG{$\p{n\sqrt{\rho}}^{\frac{1}{2\beta+1}}/\sqrt{\rho}$}
%\CL{Après vérification, je retombe sur}
$n \sqrt{\rho}$ and
$\frac{n}{\p{n \sqrt{\rho}}^{\frac{2 \beta}{2 \beta + 1}}} = \p{n\sqrt{\rho}}^{\frac{1}{2\beta+1}}/\sqrt{\rho}$
are big enough,
\begin{equation*}
    \begin{aligned}
    \inf_{\hat{\pi} \text{ $\epsilon$-DP}} \sup_{\pi \in \Theta^{\text{PSob}}_{L, \beta}}
    \E_{\vect{X} \sim \Prob_{\pi}^{\otimes n}, \hat{\pi}} \int_{[0, 1]}\p{\hat{\pi}(\vect{X})- \pi}^2
    \geq 
     C''^{-1} \p{n \sqrt{\rho}}^{-\frac{
     %2
    2\beta}{\beta + 1/2}} \;,
    \end{aligned}
\end{equation*}
%\RG{Verifier l'exposant final? Il me semble que c'est $m^{-2\beta}$ comme précédemment, ce qui donne ce que j'ai mis en rouge? \CL{Oui mais j'avais divisé le dénominateur par deux également} A corriger si nécessaire dans le texte principal et les discussions correspondantes}
where $C''$ is a constant depending only on $L$ and $\beta$.

\paragraph{Assouad version.}

From \Cref{eq:distkernelsaw}, we can see that when $h \eqdef \frac{c}{m}$ for a positive $c$ that is small enough, the condition expressed in \Cref{eq:Assouadcondition} is satisfied for $\tau = \Omega(h^{2 \beta + 1})$. To apply~\eqref{eq:AssouadMixtures}, the only missing ingredient is to bound the testing difficulties between the mixtures on the hypercube.

In the sequel, $\Prob_{\omega}$ is used as a short for $\Prob_{g_{L, \beta, \omega, h}}$. We need to bound the total variation between the mixtures on the hypercube (denoted 
$\Prob_{\omega^{i, 0}}$ and $\Prob_{\omega^{i, 1}}$, cf \eqref{eq:AssouadMixtures})
as
%\RG{Ici aussi, ne gagnerait-on pas à factoriser avec un calcul déjà fait pour se focaliser sur ce qui est spécifique à ce cas?}
%\RG{Par ailleurs, comme précédemment, $2^{-m}$ ne devrait-il pas etre $2^{-(m-1)}$? \CL{Non car la formule est $tv = 1/2 \int |f-f'|$}}
\begin{equation*}
\begin{aligned}
    %&
&\tv{\Prob_{\omega^{i, 0}}}{\Prob_{\omega^{i, 1}}}\\
%    \tv{\frac{1}{2^{m-1}} \sum_{\omega \in \{0, 1 \}^m | \omega_i = 0} \Prob_{g_{L, \beta, \omega, h}} }{\frac{1}{2^{m-1}} \sum_{\omega \in \{0, 1 \}^m | \omega_i = 1} \Prob_{g_{L, \beta, \omega, h}} } \\
    &\quad\quad\quad\quad= \frac{1}{2} \frac{1}{2^{m-1}} \int \left| \sum_{\omega \in \{0, 1 \}^m | \omega_i = 0} g_{L, \beta, \omega, h} - \sum_{\omega \in \{0, 1 \}^m | \omega_i = 1} g_{L, \beta, \omega, h} \right| \\
    &\quad\quad\quad\quad= \frac{1}{2^m} \int \Bigg| \sum_{\omega_1, \dots, \omega_{i-1}, \omega_{i+1} \dots, \omega_{m} \in \{0, 1 \}} \big( g_{L, \beta, (\omega_1, \dots, \omega_{i-1}, 0, \omega_{i+1} \dots, \omega_{m}), h} - \\  
    &\quad\quad\quad\quad\quad\quad\quad\quad\quad\quad\quad\quad\quad\quad\quad\quad\quad\quad\quad\quad\quad\quad\quad\quad g_{L, \beta, (\omega_1, \dots, \omega_{i-1}, 1, \omega_{i+1} \dots, \omega_{m}), h} \big) \Bigg| \\
    &\quad\quad\quad\quad\leq \frac{1}{2^m} \sum_{\omega_1, \dots, \omega_{i-1}, \omega_{i+1} \dots, \omega_{m} \in \{0, 1 \}} \int \Bigg|   g_{L, \beta, (\omega_1, \dots, \omega_{i-1}, 0, \omega_{i+1} \dots, \omega_{m}), h} - \\  
    &\quad\quad\quad\quad\quad\quad\quad\quad\quad\quad\quad\quad\quad\quad\quad\quad\quad\quad\quad\quad\quad\quad\quad\quad g_{L, \beta, (\omega_1, \dots, \omega_{i-1}, 1, \omega_{i+1} \dots, \omega_{m}), h}  \Bigg| \\
     &\quad\quad\quad\quad\stackrel{\eqref{eq:tvkernelsawham}}{\leq} \frac{1}{2^m} \sum_{\omega_1, \dots, \omega_{i-1}, \omega_{i+1} \dots, \omega_{m} \in \{0, 1 \}}  2  L h^{\beta + 1} \int K \\
    &\quad\quad\quad\quad= O \p{h^{\beta + 1}}\;.
\end{aligned}  
\end{equation*}
%Now, we can notice that in this sum, $(\omega_1, \dots, \omega_{i-1}, 0, \omega_{i+1} \dots, \omega_{m})$ and $(\omega_1, \dots, \omega_{i-1}, 1, \omega_{i+1} \dots, \omega_{m})$ are always at Hamming distance $1$. The difference $g_{L, \beta, (\omega_1, \dots, \omega_{i-1}, 0, \omega_{i+1} \dots, \omega_{m}), h} - g_{L, \beta, (\omega_1, \dots, \omega_{i-1}, 1, \omega_{i+1} \dots, \omega_{m}), h}$ can thus be simplified a lot. By plugging back the definition \Cref{eq:defgkernel}, and by the triangular inequality, we get

%\begin{equation*}
%\begin{aligned}
    %&\tv{\frac{1}{2^{m-1}} \sum_{\omega \in \{0, 1 \}^m | \omega_i = 0} \Prob_{g_{L, \beta, \omega, h}} }{\frac{1}{2^{m-1}} \sum_{\omega \in \{0, 1 \}^m | \omega_i = 1} \Prob_{g_{L, \beta, \omega, h}} } \\
%    \RG{\tv{\Prob_{\omega^{i, 0}}}{\Prob_{\omega^{i, 1}}}}
%    &\quad\quad\quad\quad\leq \frac{1}{2^m} \sum_{\omega_1, \dots, \omega_{i-1}, \omega_{i+1} \dots, \omega_{m} \in \{0, 1 \}} \int \Bigg(  L h^{\beta + 1} \int K + L h^\beta K \p{\frac{\cdot}{h}} \Bigg) \\
%    &\quad\quad\quad\quad\leq \frac{1}{2^m} \sum_{\omega_1, \dots, \omega_{i-1}, \omega_{i+1} \dots, \omega_{m} \in \{0, 1 \}}  2  L h^{\beta + 1} \int K \\
%    &\quad\quad\quad\quad= O \p{h^{\beta + 1}}\;.
%\end{aligned}  
%\end{equation*}

Here and in the sequel, the asymptotic comparators only hide constants (such as $\int K$ or $\int K^2$) and terms that depends on $L$ and $\beta$.
All in all, by using Le Cam's lemma %\RG{Tout ces lemmes de notre autre papiers sont utilisés si souvent que je pense qu'il faudrait les recopier dans une section au début de l'annexe afin d'éviter au lecteur de de voir aller les rechercher. Pour ma part je n'ai pas été véifier } 
for product distribution and $\rho$-zCDP (\Cref{fact:lecamslemmaconcentrated}),  and by leveraging \Cref{eq:Assouadprivate}, with $\tau = \Omega(h^{2\beta+1})$, 
\begin{equation}
    \inf_{\hat{\pi} \text{ $\rho$-zCDP}} \sup_{\pi \in  \Theta^{\text{PSob}}_{L, \beta}}
    \E_{\vect{X} \sim \Prob_{\pi}^{\otimes n}, \hat{\pi}} \int_{[0, 1]}\p{\hat{\pi}(\vect{X})- \pi}^2 = \Omega \p{m h^{2 \beta + 1}} \p{1 - n \sqrt{\rho} O \p{h^{\beta + 1}}}.
\end{equation}
%\RG{Rappeler plutot que $h = c/m$ avec $c$ suffisamment petit, et spécifier $m$ comme précédemment}
Setting $h \approx \p{n \sqrt{\rho}}^{\frac{-1}{\beta + 1}}$ and $m = c/h$ for $c$ small enough concludes the proof by yielding a lower bound $\Omega\p{\p{n\sqrt{\rho}}^{-\frac{2\beta}{2\beta+1}}}$.

\end{document}
